% ============================================================================
% STA-HPD conference paper.
%
% This file is heavily commented with "why" notes (not "what" notes) explaining
% the reasoning behind each structural and content decision. The comments are
% LaTeX comments (%) and do not affect the compiled output; they exist so that
% anyone editing this file later understands *why* a given choice was made,
% not just what the choice was.
% ============================================================================

% 10pt + twocolumn + a4paper: the user asked for a conference-style paper
% capped at roughly 3 pages of formatted density. Two-column at 10pt is the
% standard way conference proceedings (IEEE/ACM style) pack more content per
% physical page than a single-column article class would.
\documentclass[10pt,twocolumn,a4paper]{article}

% margin=0.75in: tighter than the article-class default (1in) for the same
% reason as above -- more usable width/height per page without going below a
% legible margin, matching the user's explicit page-budget constraint.
\usepackage[margin=0.75in]{geometry}
% amsmath/amssymb/amsfonts: needed for \operatorname, \mathbb, aligned
% equations, and the assorted math symbols used throughout Section II.
% amsthm: provides \newtheorem and the \begin{proof}...\end{proof}
% environment used for Theorems 1 and 2 -- the paper needs actual proofs,
% not just stated results, so amsthm's proof environment is load-bearing.
\usepackage{amsmath, amssymb, amsfonts, amsthm}
% bm: bold vectors (\bm x) for the covariate, matching the notation already
% used in this project's calibration_methods.tex so the two documents are
% typographically consistent.
\usepackage{bm}
% booktabs: the user's spec explicitly requires booktabs-style tables
% (\toprule/\midrule/\bottomrule) rather than default LaTeX table rules.
\usepackage{booktabs}
% algorithm + algpseudocode: Section II needs a concrete, citable fitting
% procedure (Algorithm 1) rather than only prose, since the tuning procedure
% (repeated splitting + mode vote) is exactly the part of the method that is
% easy to get wrong if only described in words.
\usepackage{algorithm}
\usepackage{algpseudocode}
% cite: the spec asks for \cite-based references; this package gives
% slightly better citation-list formatting/sorting than the bare kernel.
\usepackage{cite}
\usepackage{graphicx}
% hyperref: makes \label/\ref/\cite cross-references and the bibliography
% clickable -- low cost, standard for anything meant to be read as a PDF.
\usepackage{hyperref}

% Explicit \newtheorem declarations: article class has no built-in Theorem/
% Proposition/Lemma/Corollary/Remark/Definition/Example environments, so
% these must be declared before use. Definition and Example are needed by
% the new "Background: highest-density regions" subsection (Sec. II-B),
% which introduces the classical HPD-region object and a worked bimodal
% instance before any conformal-specific content is layered on top of it.
\newtheorem{theorem}{Theorem}
\newtheorem{proposition}{Proposition}
\newtheorem{lemma}{Lemma}
\newtheorem{corollary}{Corollary}
\newtheorem{remark}{Remark}
\newtheorem{definition}{Definition}
\newtheorem{example}{Example}

% Title spells out both the acronym and the application domain (MoE
% regression) up front, since a reader skimming a conference table of
% contents needs to know this is a regression-calibration paper, not a
% classification or vision paper, from the title alone.
\title{\Large STA-HPD: Shape--Threshold Adaptive Highest-Density-Region\\
Conformal Prediction for Mixture-of-Experts Regression}

% IEEEauthorblockN/A are normally supplied by the IEEEtran class. Since this
% document intentionally uses the plain `article` class (per the user's
% explicit \documentclass spec) rather than IEEEtran, the two commands are
% redefined locally just below so the familiar IEEE-style author block markup
% can still be used without pulling in the whole IEEEtran class.
\author{
\IEEEauthorblockN{Gilad Aviv}
\IEEEauthorblockA{Department of Statistics and Data Science\\
\texttt{giladaviv987@gmail.com}}
}

% See note above \author{}: these two \newcommand definitions are what make
% \IEEEauthorblockN / \IEEEauthorblockA valid commands under the plain
% article class. Kept immediately after \author so the dependency is visible
% in one place rather than buried earlier in the preamble.
\newcommand{\IEEEauthorblockN}[1]{\textbf{#1}}
\newcommand{\IEEEauthorblockA}[1]{#1}

\date{}

\begin{document}
% \twocolumn[...\begin{@twocolumnfalse}...] is the standard trick for
% getting a single-column-width title/abstract/keywords block at the top of
% an otherwise two-column document (plain \maketitle inside a twocolumn
% document would instead render the title in a single (narrow) column).
\twocolumn[
\begin{@twocolumnfalse}
\maketitle
% Abstract is capped at ~150 words per the user's spec. It is written to
% state, in order: the problem with existing methods (symmetric scaling vs.
% fixed-density HPD), the proposed fix (two-parameter reshaping/tilting),
% the theoretical guarantee (validity for any fixed (c,theta)), and the
% headline empirical number -- the order a reviewer skims an abstract in.
\begin{abstract}
Conformal prediction for heteroscedastic regression typically scales a symmetric interval by a
scalar uncertainty estimate, or, when the model exposes a full predictive density, thresholds that
density at a single global level (highest-density-region, HPD, calibration). Both choices are
suboptimal in different regimes: symmetric scaling cannot exploit multimodal expert disagreement,
while fixed-density HPD calibration inherits any systematic mis-scaling of the model's own
heteroscedasticity uncoerrected. We present STA-HPD (Shape--Threshold Adaptive HPD), which embeds
plain HPD calibration inside a two-parameter family that reshapes each mixture component's variance
by an exponent $c$ and tilts the density threshold by a scale-dependent field with exponent
$\theta$, both selected by conformal risk minimization on a held-out fold with repeated-split
mode-voting. We show the construction preserves exact finite-sample marginal coverage for any fixed
$(c,\theta)$, derive the population-level optimality condition it targets, and show it strictly
contains Standard CP, variance-scaled CP, and plain HPD calibration as boundary cases. Simulation
results show consistent width reductions of $1$--$23\%$ over plain HPD calibration at matched
coverage.
\end{abstract}
% Keywords chosen to be the terms a reader would actually search on: the
% method family (conformal prediction / distribution-free calibration), the
% application (uncertainty quantification, MoE), and the specific technique
% (HPD regions, heteroscedastic regression) that differentiates this paper
% from a generic conformal-prediction survey.
\textbf{Keywords:} conformal prediction, uncertainty quantification, mixture of experts, highest
density regions, distribution-free calibration, heteroscedastic regression
\end{@twocolumnfalse}
]

\section{Introduction}
% Why this opening: start from the most standard, least controversial fact
% (split conformal gives distribution-free coverage) so a reader unfamiliar
% with any of the specific baselines still has firm footing before the
% paper narrows to MoE-specific concerns.
Split conformal prediction \cite{vovk2005algorithmic,lei2018distribution} turns any point predictor
into a set predictor with a distribution-free, finite-sample marginal coverage guarantee, by
conformalizing a nonconformity score on a held-out calibration fold. The classical instantiation
conformalizes the raw absolute residual and returns a single interval of constant width. Follow-up
work makes the interval heteroscedastic by normalizing the residual with a model-reported scale
(variance-scaled conformal prediction) \cite{lei2018distribution}, or by conformalizing a
quantile-regression head directly (CQR) \cite{romano2019conformalized}. Both remain committed to a
symmetric, single-interval region $\hat y(\bm x) \pm s(\bm x)$.

% Why this paragraph: this is the actual motivating gap the paper closes --
% without it, a reader could reasonably ask "why not just use CQR or CP-VS
% everywhere?". Making the MoE-multimodality argument explicit, with the
% Neyman-Pearson pointer forward to Theorem 1, sets up why HPD calibration
% (and by extension STA-HPD) is the right tool for this model class
% specifically, not conformal prediction in general.
This commitment is a real cost when the underlying model is a Mixture-of-Experts (MoE) regressor
\cite{jordan1994hierarchical}.
An MoE model with a probabilistic head naturally represents its predictive uncertainty as a Gaussian
mixture density over the $K$ expert components, and this density is routinely multimodal: distinct
experts can disagree not about the scale of the outcome, but about which of several plausible
outcomes is correct. A symmetric interval built around the aggregate mean is provably wasteful in
this regime, since it must span the full range between well-separated modes while covering the
low-density valley between them. Highest-density-region (HPD) conformal calibration
\cite{izbicki2020cd,sadinle2019least} avoids this by conformalizing the negative log-density and
returning the induced super-level set of the fitted mixture, which is width-optimal \emph{for the
fitted density itself} by a Neyman--Pearson argument.

% Why this paragraph: plain HPD calibration is not itself the paper's
% contribution, so before introducing STA-HPD the paper has to state
% precisely what is left broken after HPD calibration is applied. The two
% numbered failure modes here map 1:1 onto the two tunable knobs (c and
% theta) introduced later -- this paragraph is effectively the specification
% for what those two knobs need to fix.
The gap this leaves is: the fitted mixture is not the true conditional density, and plain HPD
calibration has no mechanism to correct \emph{how} that mismatch is structured. Two systematic
error modes recur in practice: (i) the model's per-component variances $\sigma_k(\bm x)$ may be
globally over- or under-dispersed relative to the residual scale, and (ii) that dispersion may be
more or less \emph{informative} about the residual magnitude on some datasets than others, so that
even a perfectly-centered variance estimate should not be trusted with a fixed, dataset-independent
curvature.

% Why a numbered contributions paragraph: reviewers scan for this
% specifically, and numbering it forces the four claims to be independently
% checkable against the rest of the paper (shape exponent -> Sec. II-C;
% threshold tilt -> Sec. II-C; tuning procedure -> Sec. II-E; boundary
% containment -> Proposition 1) rather than left as a vague summary.
\paragraph{Contributions.} We introduce STA-HPD, a two-parameter generalization of HPD calibration
that (1) reshapes per-component variances by a power exponent $c$ before recomputing the mixture
density, correcting systematic over/under-dispersion; (2) tilts the conformal threshold itself by a
power of the point's total predictive scale, exponent $\theta$, correcting how strongly that scale
should move the threshold; (3) selects $(c,\theta)$ by conformal risk minimization under a
repeated-splitting, mode-voting procedure that keeps the search itself statistically valid; and (4)
provably contains Standard CP, variance-scaled CP (in the single-expert case), and plain HPD
calibration as specific grid points, so the tuned method's own tuning objective can never be worse
than the strongest of its ancestors on the tuning fold.

% Why a dedicated Related Work section (rather than folding these citations
% into the Introduction): two of the four lines below -- post-hoc score
% *reshaping* for width (Boosted CP) and explicit length *minimization*
% subject to validity (CPL) -- share STA-HPD's high-level objective closely
% enough that a reviewer familiar with the 2024 conformal literature would
% immediately ask how STA-HPD differs. Answering that head-on, in its own
% section, is what lets the Contributions paragraph above claim a *narrow,
% specific* novelty (the per-component shape exponent + the closed-form
% nested family) honestly, rather than the broader "post-hoc width tuning"
% claim which is no longer new. Each paragraph names the closest prior work
% and states precisely what STA-HPD does that it does not.
\section{Related Work}
\label{sec:related}

\paragraph{Density-based and highest-density conformal regions.} The idea of conformalizing a
predictive density and returning its super-level set, rather than a symmetric interval, is due to
Izbicki et al.\ (CD-split / HPD-split) \cite{izbicki2020cd} and, for the set-valued classification
analogue, Sadinle et al.\ \cite{sadinle2019least}; both build on the classical highest-density-region
object \cite{boxtiao1973bayesian,hyndman1996computing}. STA-HPD's $(c,\theta)=(1,0)$ boundary case
(plain MoG-HPD) is exactly a mixture-density instance of this construction and is \emph{not} claimed
as novel here. Our contribution begins one level up: rather than leveling the fitted density as
given, we level a \emph{reshaped} density $f_c$ and tilt its threshold, and show (Theorem~\ref{thm:optimality})
why the raw fitted density is generally the wrong one to level when it is only an estimate of the
true conditional law.

\paragraph{Normalized / locally adaptive scores.} Scaling the residual by a reported predictive
scale, $|y-\hat y|/\sigma(\bm x)$ \cite{lei2018distribution,papadopoulos2011regression}, and
conformalizing a learned quantile head (CQR) \cite{romano2019conformalized}, are the standard routes
to heteroscedastic width. Both hard-code how width scales with the reported uncertainty: CP-VS fixes
width $\propto\sigma(\bm x)$, i.e.\ the elasticity $\partial\log s/\partial\log\sigma$ is pinned at
$1$. STA-HPD's shape exponent $c$ is exactly a free such elasticity (Lemma~\ref{lem:logaffine}),
and its threshold field $\theta$ recovers the CP-VS scaling as the single point $(1,1)$
(Corollary~\ref{cor:k1closedform}); the empirical finding that the width-minimizing $c^\star$ is
often strictly below $1$ is precisely a statement that this elasticity should not be fixed at the
CP-VS value.

\paragraph{Post-hoc width minimization.} Closest in spirit is Boosted Conformal Prediction
\cite{xie2024boosted}, which is likewise post-hoc (operating on a frozen model) and likewise
\emph{reshapes a base conformity score specifically to reduce interval length at valid coverage} --
so the high-level strategy of ``tune the score after training to cut width'' is not itself new to
this paper. STA-HPD differs in mechanism and scope: Boosted CP learns a flexible score correction by
gradient boosting over contrast trees, whereas STA-HPD uses a closed-form two-parameter family whose
members are named calibrators (Proposition~\ref{prop:recovery}), selected over a small discrete grid
by mode-voting -- yielding an interpretable optimum $(c^\star,\theta^\star)$ and a mixture-specific,
multimodal-capable region shape (Eq.~\eqref{eq:region}) that a tree-boosted scalar score does not
target. Length-optimization approaches that minimize expected set size subject to (conditional)
validity via constrained optimization or duality \cite{kiyani2024length} share STA-HPD's objective
but operate as general-purpose optimization frameworks rather than a mixture-density family with
closed-form boundary cases. STA-HPD should be read as a specialized, low-degree-of-freedom point in
this design space, chosen for interpretability and for exactly containing the MoE calibrators it
must beat, not as a competitor to the full generality of those methods.

\paragraph{Stabilizing data-dependent selection.} The repeated-split, mode-vote / median-quantile
procedure (Sec.~\ref{sec:math}-E) is a cross-conformal-style aggregation in the lineage of
\cite{gibbs2021adaptive,tibshirani2019conformal}; here it serves only to keep the discrete
$(c,\theta)$ search from being dominated by the noise of any single tune/calib split, and is not a
contribution in its own right.

\section{Mathematical Foundations of STA-HPD}
\label{sec:math}

\subsection{Notation and setup}
% Why notation is introduced before any theorem: Theorem 1 needs "a density
% f" only, but everything from Theorem 2 onward is stated in terms of the
% MoE's specific (w_k, mu_k, sigma_k) outputs, so those symbols have to be
% fixed once, here, rather than redefined ad hoc in each subsection.
Let $\mathcal D_{\mathrm{cal}} = \{(\bm x_i, y_i)\}_{i=1}^n$ be an exchangeable calibration sample
and let a trained $K$-expert probabilistic MoE model supply, for any $\bm x$, gating weights
$w_k(\bm x)$ with $\sum_k w_k(\bm x)=1$, per-expert means $\mu_k(\bm x)$, and per-expert predictive
standard deviations $\sigma_k(\bm x)$. The aggregate point prediction is
$\hat y(\bm x) = \sum_{k=1}^K w_k(\bm x)\,\mu_k(\bm x)$ and the model's exact mixture predictive
density is
\begin{equation}
f_{\mathrm{mix}}(y\mid \bm x) = \sum_{k=1}^{K} w_k(\bm x)\, \mathcal N\big(y;\ \mu_k(\bm x),\ \sigma_k^2(\bm x)\big).
\label{eq:mix}
\end{equation}
% Why q_level gets its own equation and label: every method in Sections
% II-III shares this exact quantile correction. Stating it once here, with a
% label, lets every later construction just say "quantile at q_level" and
% cite eq:qlevel, instead of re-deriving the finite-sample correction each
% time -- and it is the single fact that makes every coverage claim in this
% paper (Prop. 2, Sec. III, Sec. IV) true regardless of model correctness.
Every calibrator in this family targets a user-specified miscoverage level $\alpha\in(0,1)$ and
uses the same finite-sample-corrected empirical quantile at level
\begin{equation}
q_{\mathrm{level}} = \min\!\left(1,\ \frac{\lceil (n+1)(1-\alpha)\rceil}{n}\right),
\label{eq:qlevel}
\end{equation}
which is the quantity that gives split conformal prediction its distribution-free marginal
coverage guarantee $P(y_{\mathrm{test}}\in\hat C(\bm x_{\mathrm{test}})) \ge 1-\alpha$ for any fixed,
data-independent nonconformity score, regardless of whether the base model is correctly specified.

\subsection{Background: highest-density regions}
\label{sec:hpdbasics}
% Why this subsection is inserted before any conformal-specific content:
% "HPD region" is a classical object from Bayesian and frequentist density
% estimation, entirely independent of conformal prediction -- Theorem 1
% below is a fact about *any* density on the real line, proved without ever
% mentioning a calibration set, a model, or a p-value. Presenting it here as
% background, before Section II-C introduces the conformal machinery that
% consumes it, makes clear that the width-optimality property STA-HPD relies
% on is a statement about level sets in general, and would be true even if
% conformal prediction had never been invented.
Before specializing to mixture-of-experts calibration, we recall the classical notion of a
highest-density region (HDR/HPD region), originally studied as the natural generalization of a
symmetric credible interval to arbitrary, possibly asymmetric or multimodal densities
\cite{boxtiao1973bayesian,hyndman1996computing}.

\begin{definition}[Highest-density region]
\label{def:hpd}
For a density $f$ on $\mathbb R$ and $\alpha\in(0,1)$, the $100(1-\alpha)\%$ highest-density region
of $f$ is
\begin{equation}
R_\alpha(f) = \{y \in \mathbb R : f(y) \ge \tau_\alpha\},
\end{equation}
where $\tau_\alpha$ is the largest constant such that $\int_{R_\alpha(f)} f \ge 1-\alpha$.
\end{definition}

% Why Theorem 1 is stated here, about a bare density f, rather than later
% about f_mix specifically: keeping it fully general is what lets Section
% II-C (MoG-HPD), the STA-HPD family (which levels a *different* density
% f_c), and even the K=1 Gaussian corollary below all invoke the *same*
% theorem without re-proving it -- the theorem does not care what f is, only
% that it is a fixed density.
\begin{theorem}[Level-set width optimality]
\label{thm:hpd}
Fix a density $f$ and $\alpha \in (0,1)$, and let $R_\tau = \{y : f(y)\ge \tau\}$ with $\tau$ chosen so
that $\int_{R_\tau} f = 1-\alpha$ (i.e.\ $R_\tau = R_\alpha(f)$ of Definition~\ref{def:hpd}). Then for
any measurable $C$ with $\int_C f \ge 1-\alpha$, $|R_\tau| \le |C|$, where $|\cdot|$ denotes Lebesgue
measure.
\end{theorem}
% Proof strategy: split the symmetric difference of R_tau and C into D_1,
% D_2 and bound each side using the fact that f is, by definition of R_tau,
% above tau on D_1 and below tau on D_2 -- the standard Neyman-Pearson
% "trade high-density mass for low-density mass" argument, spelled out so
% the reader can verify each inequality rather than take the conclusion on
% faith.
\begin{proof}
Let $D_1 = R_\tau\setminus C$ and $D_2 = C\setminus R_\tau$. Since $D_1\subset R_\tau$, $f\ge\tau$ on
$D_1$; since $D_2\subset R_\tau^c$, $f<\tau$ on $D_2$. As $R_\tau\cap C$ is common to both sets,
$\int_C f\ge\int_{R_\tau}f$ forces $\int_{D_2}f\ge\int_{D_1}f$. Hence
$\tau|D_2|\ge\int_{D_2}f\ge\int_{D_1}f\ge\tau|D_1|$, so $|D_2|\ge|D_1|$ and
$|R_\tau|-|C| = |D_1|-|D_2|\le 0$.
\end{proof}

% Why an explicit Gaussian corollary is worth stating on its own: it is the
% bridge back to every symmetric-interval baseline in Section III (Standard
% CP, CP-VS). Without it, a reader could reasonably wonder how a *level-set*
% construction relates to the *symmetric-interval* constructions used
% everywhere else in the conformal-prediction literature; this corollary
% shows the two coincide exactly in the one case (unimodal Gaussian) where
% symmetric intervals are actually correct.
\begin{corollary}[Unimodal Gaussian reduction]
\label{cor:gaussian-hpd}
If $f = \mathcal N(\mu,\sigma^2)$, then $R_\alpha(f)$ is the single symmetric interval
$[\mu - \sigma z_{1-\alpha/2},\ \mu + \sigma z_{1-\alpha/2}]$, where $z_{1-\alpha/2}$ is the
$(1-\alpha/2)$ standard-normal quantile.
\end{corollary}
\begin{proof}
$f$ is symmetric and strictly decreasing in $|y-\mu|$, so every super-level set $\{f\ge\tau\}$ is an
interval centered at $\mu$; choosing its half-width so that it carries mass $1-\alpha$ under a
standard Gaussian gives exactly $\sigma z_{1-\alpha/2}$.
\end{proof}

% Why this remark on equal-tailed regions is included: it names the
% *implicit* assumption behind every symmetric-interval baseline in this
% paper (Standard CP, CP-VS) explicitly, rather than leaving it implicit --
% those methods are only shape-optimal when the residual distribution
% happens to be unimodal and symmetric, which is exactly the boundary case
% Corollary~\ref{cor:gaussian-hpd} identifies.
\begin{remark}[Equal-tailed vs.\ highest-density regions]
The equal-tailed region $[F^{-1}(\alpha/2), F^{-1}(1-\alpha/2)]$, with $F$ the CDF of $f$, is the
construction implicitly targeted by every symmetric-interval calibrator in this paper (Standard CP,
CP-VS). By Theorem~\ref{thm:hpd} it can only be as narrow as $R_\alpha(f)$, with equality when $f$ is
unimodal and symmetric (Corollary~\ref{cor:gaussian-hpd}); for skewed or multimodal $f$ the
equal-tailed region is strictly wider, since it fixes \emph{where} probability mass is excluded
(equally in each tail) rather than \emph{how much} density is required to be included.
\end{remark}

% Why a fully worked bimodal example, with concrete numbers, rather than
% just asserting "HPD regions can be disjoint": an abstract claim about
% multimodality is easy to accept and forget; a numeric instance the reader
% can (in principle) re-derive by hand is what makes concrete, later, why
% Section IV's Synthetic dataset -- whose true law is a genuine mixture --
% is exactly where STA-HPD's advantage over symmetric methods is largest.
\begin{example}[Disjoint highest-density regions]
\label{ex:bimodal}
Let $f = 0.5\,\mathcal N(2,0.25) + 0.5\,\mathcal N(8,0.25)$, a two-component mixture with
well-separated means and small component variances. The aggregate mean is $\hat y = 5$, but
$f(5)$ is lower than $f$ near either mode, since $5$ sits in the low-density valley between them. If
$\tau_\alpha$ falls between $f(5)$ and the peak heights, $R_\alpha(f)$ is two disjoint intervals
bracketing $2$ and $8$ and \emph{excludes} the region around $\hat y=5$ entirely -- whereas any
symmetric interval $\hat y\pm s$ must either waste width covering that empty middle region or fail
to reach one of the two modes.
\end{example}

\subsection{Conformalizing HPD regions: MoG-HPD}
% Why this paragraph follows immediately: Theorem 1 is abstract (any density
% f); this paragraph is what actually turns it into a usable calibrator by
% plugging in f_mix and the negative-log-density score, and states clearly
% that validity (from eq:qlevel) and optimality (from Theorem 1) are two
% separate claims -- one holds regardless of model quality, the other only
% holds relative to whatever density is being leveled.
Plain HPD calibration (MoG-HPD) applies Theorem~\ref{thm:hpd} with $f = f_{\mathrm{mix}}(\cdot\mid\bm x)$
by conformalizing the negative log-density,
\begin{equation}
S_i = -\log f_{\mathrm{mix}}(y_i\mid\bm x_i), \qquad \hat t = \operatorname{Quantile}_{q_{\mathrm{level}}}(\{S_i\}),
\end{equation}
and returning $\hat C(\bm x) = \{y: f_{\mathrm{mix}}(y\mid\bm x)\ge e^{-\hat t}\}$, i.e.\ the empirical
$R_{\alpha}(f_{\mathrm{mix}}(\cdot\mid\bm x))$ of Definition~\ref{def:hpd} at the conformalized level
$\hat\tau=e^{-\hat t}$. Validity follows from Eq.~\eqref{eq:qlevel} regardless of whether
$f_{\mathrm{mix}}$ is well specified; only the achieved width depends on how close $f_{\mathrm{mix}}$
is to the true conditional density -- the same validity/efficiency split noted for
Example~\ref{ex:bimodal}, now under sampling uncertainty in $\hat\tau$ rather than a known $\tau$.

\subsection{The STA-HPD family}
% Why introduce G (the frozen geometric mean) before the shape exponent:
% c needs a reference scale to exponentiate around, or c != 1 would shift
% the *absolute* scale of every component instead of just its *spread*
% relative to the calibration-set typical scale. Freezing G from the
% calibration set (not per test point) keeps the reshaping a fixed,
% label-independent transform, which Proposition 2's validity argument later
% depends on.
Plain HPD calibration fixes the density used to measure width to be exactly $f_{\mathrm{mix}}$, and
thresholds it at a single point-independent level. STA-HPD opens both choices. Freeze the
calibration-set geometric mean of per-component scales,
\begin{equation}
G = \exp\!\Big(\operatorname*{mean}_{k,i}\ \log\sigma_k(\bm x_i)\Big),
\end{equation}
and define, for a \emph{shape} exponent $c\in\mathbb R_{\ge0}$, the reshaped component scale and the
corresponding surrogate mixture density
\begin{align}
\tilde\sigma_k(\bm x; c) &= G\,\big(\sigma_k(\bm x)/G\big)^{c}, \label{eq:shape}\\
f_c(y\mid\bm x) &= \sum_{k=1}^K w_k(\bm x)\,\mathcal N\big(y;\ \mu_k(\bm x),\ \tilde\sigma_k^2(\bm x;c)\big). \label{eq:fc}
\end{align}
The exponent acts multiplicatively in log-scale space: $c<1$ compresses the spread of component
scales toward the calibration-set typical scale $G$ (shrinking heteroscedasticity), $c>1$ amplifies
it, and $c=1$ recovers $f_{\mathrm{mix}}$ exactly. The next lemma makes this precise.

% Why this is stated as a formal Lemma rather than left as the prose
% sentence above: "acts multiplicatively in log-scale space" is a claim that
% is trivial to verify but easy to get subtly wrong when reasoning about it
% informally (e.g. mistaking it for a claim about the variance rather than
% the standard deviation, or about an additive rather than affine
% transform). Writing it out removes any ambiguity about exactly what
% "compresses toward G" means, and the three named consequences are exactly
% the three boundary values ($c=1,0,\to\infty$) that Proposition 1 and
% Table I later depend on.
\begin{lemma}[Shape reshaping is affine in log-scale]
\label{lem:logaffine}
For every $k$ and $\bm x$, $\log\tilde\sigma_k(\bm x;c) = \log G + c\big(\log\sigma_k(\bm x) - \log G\big)$.
Consequently $\tilde\sigma_k(\bm x;1)=\sigma_k(\bm x)$ for all $k,\bm x$; $\tilde\sigma_k(\bm x;0)=G$
for all $k,\bm x$ (every component collapses to the same constant scale); and for $c>0$,
$\tilde\sigma_k(\bm x;c)$ is monotonically increasing in $\sigma_k(\bm x)$ with elasticity exactly
$c$, i.e.\ $\partial\log\tilde\sigma_k/\partial\log\sigma_k = c$.
\end{lemma}
\begin{proof}
Take logarithms of Eq.~\eqref{eq:shape}: $\log\tilde\sigma_k(\bm x;c) = \log G + c\log(\sigma_k(\bm x)/G)
= \log G + c(\log\sigma_k(\bm x)-\log G)$, which is affine in $\log\sigma_k(\bm x)$ with slope $c$ and
value $\log G$ at $\log\sigma_k(\bm x)=\log G$. Substituting $c=1$ and $c=0$ gives the two stated
boundary identities; differentiating the affine form with respect to $\log\sigma_k(\bm x)$ gives the
elasticity $c$ directly.
\end{proof}

Define the point's own total predictive scale
under the reshaped family,
\begin{equation}
\sigma_{\mathrm{tot}}^2(\bm x; c) = \sum_{k=1}^K w_k(\bm x)\,\tilde\sigma_k^2(\bm x;c) \;+\; \sum_{k=1}^K w_k(\bm x)\,\big(\mu_k(\bm x)-\hat y(\bm x)\big)^2,
\end{equation}
% Why sigma_tot uses the reshaped aleatoric term but an untouched epistemic
% term: c is defined (above) to act only on each component's own variance
% sigma_k, not on how far apart the component means are, so the epistemic
% (between-expert) part of the variance decomposition is left exactly as
% the base model reported it -- c and theta are meant to correct the
% *aleatoric* scale specifically, per the two failure modes named in the
% Introduction.
the reshaped aleatoric term plus the (shape-invariant) epistemic term. The \emph{threshold-field}
exponent $\theta\in\mathbb R$ then tilts the conformal score by this point-wise scale,
\begin{equation}
S_i(c,\theta) = -\log f_c(y_i\mid\bm x_i) \;-\; \theta\,\log\sigma_{\mathrm{tot}}(\bm x_i; c),
\label{eq:score}
\end{equation}
with threshold $\hat t(c,\theta) = \operatorname{Quantile}_{q_{\mathrm{level}}}(\{S_i(c,\theta)\})$
and induced prediction set
\begin{equation}
\hat C(\bm x) = \Big\{ y :\ \log f_c(y\mid\bm x) \;\ge\; -\hat t(c,\theta) - \theta\log\sigma_{\mathrm{tot}}(\bm x;c) \Big\}.
\label{eq:region}
\end{equation}
% Why this sentence is here: it is easy to misread eq:region as "a
% different density thresholded globally" -- this sentence heads that off
% by naming exactly which part is still global (\hat t) and which part now
% varies by point (the sigma_tot^{-theta} tilt), which is the whole point of
% adding theta on top of c.
Equation~\eqref{eq:region} is exactly a level set of $f_c$, but at a threshold that is scaled up or
down per test point by its own $\sigma_{\mathrm{tot}}(\bm x;c)^{-\theta}$, rather than at the single
global level plain HPD uses.

% Why Proposition 1 exists: it is the formal version of the "contains
% Standard CP / CP-VS / MoG-HPD as boundary cases" claim made in the
% Contributions paragraph and in the abstract -- stating it as a numbered
% proposition makes it something Section III and IV can cite by name
% (e.g. "as Proposition 1 guarantees") rather than re-asserting informally
% each time it is used.
\begin{proposition}[Boundary recovery]
\label{prop:recovery}
$(c,\theta)=(1,0)$ recovers plain HPD calibration exactly. When $K=1$ (a single Gaussian
predictive component), $(c,\theta)=(1,1)$ recovers variance-scaled conformal prediction (CP-VS) on
the total predictive standard deviation exactly. $(c,\theta)=(0,0)$ collapses every
$\tilde\sigma_k(\bm x;0)=G$ to a shared constant, making $f_0(\cdot\mid\bm x)$ depend on $\bm x$ only
through $(\mu_k(\bm x), w_k(\bm x))$ at fixed scale, and $\theta=0$ leaves the threshold untilted,
recovering Standard (unscaled) split conformal prediction on the aggregate residual in the
single-component reduction.
\end{proposition}

% Why the proof is a short paragraph, not a formal \begin{proof} block: each
% of the three cases in Proposition 1 is direct substitution into equations
% already stated above (eq:shape through eq:region), so a full proof
% environment would only repeat those equations without adding content --
% a pointer to which substitution does the work is sufficient here.
The proof is immediate substitution into Eqs.~\eqref{eq:shape}--\eqref{eq:region}: at $c=1$,
$\tilde\sigma_k \equiv \sigma_k$ so $f_1=f_{\mathrm{mix}}$, and $\theta=0$ removes the tilt term from
Eq.~\eqref{eq:score}, giving back plain HPD's score exactly.

% Why this closed-form corollary is added on top of the substitution proof
% above: substitution shows the *scores* coincide, but does not, by itself,
% show what the resulting *region* looks like -- a reader could still
% wonder whether "the score matches CP-VS's score" really implies "the
% region is CP-VS's symmetric interval." Deriving the region explicitly for
% $K=1$ removes that gap and gives a genuinely new, checkable identity
% (the closed-form half-width $s(\bm x)$) rather than repeating the earlier
% substitution argument in different words.
\begin{corollary}[Closed-form half-width, single component]
\label{cor:k1closedform}
When $K=1$, write $\tilde\sigma(\bm x;c)$ for the (only) reshaped component scale and
$\hat\tau=e^{-\hat t(c,\theta)}$. Then $\hat C(\bm x)$ of Eq.~\eqref{eq:region} is the symmetric
interval $\hat y(\bm x)\pm s(\bm x)$ with
\begin{equation}
s(\bm x) = \tilde\sigma(\bm x;c)\ \sqrt{2\log\!\left(\frac{1}{\hat\tau\sqrt{2\pi}\ \tilde\sigma(\bm x;c)^{1-\theta}}\right)}.
\label{eq:k1closedform}
\end{equation}
In particular: at $(c,\theta)=(1,0)$, $\tilde\sigma=\sigma$ and $s(\bm x)=\sigma(\bm x)\sqrt{2\log(1/(\hat\tau\sqrt{2\pi}\,\sigma(\bm x)))}$,
plain MoG-HPD's Gaussian half-width; at $(c,\theta)=(1,1)$, the $\sigma(\bm x)^{1-\theta}$ factor
vanishes and $s(\bm x)=\sigma(\bm x)\cdot\sqrt{2\log(1/(\hat\tau\sqrt{2\pi}))}$ is exactly CP-VS's
form $\hat q\,\sigma(\bm x)$ for a constant $\hat q$; and at $(c,\theta)=(0,0)$,
$\tilde\sigma\equiv G$ is constant in $\bm x$, so $s(\bm x)\equiv G\sqrt{2\log(1/(\hat\tau\sqrt{2\pi}G))}$
is Standard CP's constant half-width.
\end{corollary}
\begin{proof}
For $K=1$, $f_c(y\mid\bm x)=\mathcal N(y;\mu(\bm x),\tilde\sigma^2(\bm x;c))$ and
$\sigma_{\mathrm{tot}}(\bm x;c)=\tilde\sigma(\bm x;c)$ exactly (the epistemic term vanishes with a
single component). Substituting the Gaussian density into $\log f_c(y\mid\bm x)\ge-\hat t-\theta\log\tilde\sigma(\bm x;c)$
and solving for $|y-\mu(\bm x)|$ gives
\[
(y-\mu(\bm x))^2 \le 2\tilde\sigma^2(\bm x;c)\Big[\hat t+\theta\log\tilde\sigma(\bm x;c)-\log\big(\tilde\sigma(\bm x;c)\sqrt{2\pi}\big)\Big],
\]
and collecting the $\log\tilde\sigma(\bm x;c)$ terms on the right-hand side, substituting
$\hat\tau=e^{-\hat t}$, and taking the square root gives Eq.~\eqref{eq:k1closedform}. The three
boundary cases follow by direct substitution of $(c,\theta)$ and Lemma~\ref{lem:logaffine}'s
identities $\tilde\sigma(\bm x;1)=\sigma(\bm x)$, $\tilde\sigma(\bm x;0)=G$.
\end{proof}

\subsection{Population-level optimality condition}
% Why this subsection exists at all: Theorem 1 only says the level set of a
% *given* density is width-optimal for *that* density -- it says nothing
% about which density to level in the first place. This subsection derives
% what the *true* population optimum would require (eq:optcond), so the gap
% between that condition and what plain HPD actually enforces is exactly
% what motivates introducing c and theta -- it is the theoretical
% justification for the method, not just an empirical after-the-fact story.
Coverage pins exactly one scalar functional of $\hat C$ (its probability mass); the two knobs
$(c,\theta)$ reallocate a fixed coverage budget across $\bm x$ without touching that scalar. Write
the per-point log-threshold as $g(\bm x)$, so that $\hat C(\bm x) = \{y: \log f(y\mid\bm x)\ge -g(\bm x)\}$
for whatever density $f$ is in force, with boundary points $b_j(\bm x)$ solving
$f(b_j(\bm x)\mid\bm x) = e^{-g(\bm x)}$.

\begin{theorem}[First-order optimality of the level-set family]
\label{thm:optimality}
Consider $\min_g \; \mathbb E_{\bm x}\big[|C_g(X)|\big]$ subject to
$\mathbb E_{\bm x}\big[P(Y\in C_g(X)\mid X=\bm x)\big] \ge 1-\alpha$, where the true conditional law
of $Y\mid \bm x$ has density $p(\cdot\mid\bm x)$. At a stationary point of the associated Lagrangian,
\begin{equation}
p\big(b(\bm x)\mid \bm x\big) = \text{const across } \bm x.
\label{eq:optcond}
\end{equation}
\end{theorem}
% Why this is now written as a full Proof (upgraded from the "Proof sketch"
% label used in the first draft) rather than left informal: the user asked
% for the theoretical explanation to be expanded, and the perturbation
% argument below is fully elementary (single real-valued perturbation
% parameter delta, first-order Taylor expansion, a single stationarity
% condition) -- there is no missing regularity assumption that would make
% "sketch" the honest label once each step is written out completely, as
% done here (each of the four numbered steps is a complete derivation, not
% a gap to be filled in by the reader).
\begin{proof}
The argument has four steps.

\emph{Step 1 (boundary perturbation).} Fix $\bm x$ and perturb the log-threshold at that point only,
$g(\bm x)\mapsto g(\bm x)+\delta$. Each boundary point $b_j(\bm x)$ solves
$\log f(b_j(\bm x)\mid\bm x)=-g(\bm x)$; implicitly differentiating in $\delta$,
$f'(b_j(\bm x))\cdot \partial b_j/\partial\delta = -f(b_j(\bm x)\mid\bm x)\cdot(-1) = f(b_j(\bm x)\mid \bm x)$
is not needed directly -- instead, linearizing $\log f(b_j+\Delta b_j\mid\bm x)\approx \log f(b_j\mid\bm x) + \Delta b_j\, f'(b_j\mid\bm x)/f(b_j\mid\bm x)$
against the shifted constraint $-g(\bm x)-\delta$ gives $\Delta b_j = -\delta\, f(b_j\mid\bm x)/f'(b_j\mid\bm x)$
to first order, i.e.\ each boundary moves by an amount inversely proportional to the density's slope
there.

\emph{Step 2 (width response).} Enlarging the region at $b_j(\bm x)$ by $|\Delta b_j|$ changes
$|C_g(\bm x)|$ by $\sum_j |\Delta b_j| = \delta\sum_j f(b_j\mid\bm x)/|f'(b_j\mid\bm x)|$ to first
order (signs chosen so increasing $\delta$, i.e.\ a lower threshold, enlarges the region).

\emph{Step 3 (coverage response).} The true covered probability
$P(Y\in C_g(X)\mid X=\bm x)=\int_{C_g(\bm x)}p(y\mid\bm x)\,dy$ changes, by the same boundary
displacement, by $\delta\sum_j p(b_j(\bm x)\mid\bm x)\,f(b_j\mid\bm x)/|f'(b_j\mid\bm x)|$ to first
order -- identical to Step 2's expression except that the model density $f(b_j\mid\bm x)$ in the
numerator is replaced by the true density $p(b_j\mid\bm x)$.

\emph{Step 4 (stationarity).} The Lagrangian is
$\mathcal L[g]=\mathbb E_{\bm x}[|C_g(X)|]-\lambda\big(\mathbb E_{\bm x}[P(Y\in C_g(X)\mid X)]-(1-\alpha)\big)$.
A stationary point requires the Fr\'echet derivative of $\mathcal L$ with respect to a pointwise
perturbation at every $\bm x$ to vanish, i.e.\ (Step 2's rate) $=\lambda\cdot$(Step 3's rate) for
every $\bm x$. Both rates share the common factor $\sum_j f(b_j(\bm x)\mid\bm x)/|f'(b_j(\bm x)\mid\bm x)|$,
which cancels, leaving $1=\lambda\,p(b(\bm x)\mid\bm x)$ (in the single-boundary case; the
multi-boundary case is a weighted average of the $p(b_j(\bm x)\mid\bm x)$ under the same weights in
both Steps 2 and 3, so the cancellation still forces the weighted-average $p$ to equal $1/\lambda$
at every $\bm x$ for a common $\lambda$). Since $\lambda$ is a single global multiplier shared across
all $\bm x$, $p(b(\bm x)\mid\bm x)=1/\lambda$ is the same constant for every $\bm x$.
\end{proof}

% Why this remark follows the proof: Theorem 2 states an abstract
% population condition; Corollary 2's closed form (eq:k1closedform) is the
% one place in the paper where "how would c and theta actually move the
% boundary" can be computed explicitly rather than left implicit in f_c and
% theta's definitions. Connecting the two turns Theorem 2 from a pure
% existence statement into something with a directly computable
% single-component illustration.
\begin{remark}[Theorem~\ref{thm:optimality} in the $K=1$ closed form]
Differentiating Eq.~\eqref{eq:k1closedform} with respect to $\tilde\sigma(\bm x;c)$ shows
$\partial s/\partial\tilde\sigma$ depends on $\theta$ only through the constant $(1-\theta)$
multiplying $\log\tilde\sigma$ inside the square root: $\theta$ rescales how fast the half-width
grows as the point's own reshaped scale grows, which is exactly the pointwise reallocation
Theorem~\ref{thm:optimality} describes -- $\theta<1$ makes $s(\bm x)$ grow faster than the shape-only
($\theta=0$) baseline in high-$\tilde\sigma$ regions and slower in low-$\tilde\sigma$ regions
(and conversely for $\theta>1$), shifting width toward or away from wherever $\tilde\sigma(\bm x;c)$
is largest without changing the marginal coverage constraint, which is fixed by $\hat t$ alone.
\end{remark}

% Why this paragraph is the payoff of the whole subsection: it explicitly
% connects eq:optcond (population truth) back to what plain HPD enforces
% (a constant *model*-density condition) and assigns each of the two new
% parameters to the specific part of that mismatch it is meant to fix --
% c to the systematic-scaling part, theta to the residual pointwise part.
% Without this paragraph, Theorem 2 would read as an interesting but
% disconnected fact rather than the method's actual design rationale.
Plain HPD calibration instead enforces $f_{\mathrm{mix}}(b(\bm x)\mid\bm x) = \hat\tau$, a constant
\emph{model}-density condition, which coincides with Eq.~\eqref{eq:optcond} only where
$f_{\mathrm{mix}}$ agrees with the true $p$ at the boundary. The shape exponent $c$ corrects the part
of this mismatch driven by systematic mis-scaling of $\sigma_k(\bm x)$ relative to residual
magnitude; the threshold field $\theta$ corrects the remaining pointwise level mismatch. In
particular, when $\sigma_k(\bm x)$ is only partially informative about $|y-\hat y(\bm x)|$, dispersion
in $\sigma_k$ inflates $\mathbb E[\sigma_k^{c}]$-type terms faster than it deflates the corresponding
conformal quantile, so the width-minimizing $c^\star$ is generally interior and can fall strictly
below $1$ — a data-driven shrinkage of heteroscedasticity unavailable to plain HPD, which fixes
$c\equiv1$ by construction.

\subsection{Fitting procedure and its validity}
% Why the tune/calib split is introduced with this specific justification
% (exchangeability) rather than just stated as a design choice: this is the
% single most common way a "tune a hyperparameter, then conformalize" method
% silently loses its coverage guarantee (by re-using labels), so the paper
% is explicit that the split is not a minor implementation detail but the
% mechanism that keeps Proposition 2 true.
Selecting $(c,\theta)$ on the same points used to fit $\hat t$ breaks exchangeability, so the
calibration set is split $\mathcal D_{\mathrm{cal}} = \mathcal D_{\mathrm{tune}}\sqcup\mathcal D_{\mathrm{calib}}$,
with $|\mathcal D_{\mathrm{tune}}| = \texttt{tune\_frac}\cdot n$ (default $0.2$, deliberately smaller
than a $50/50$ split since the grid is low-dimensional and discrete while $\hat t$ is a continuous
quantile that benefits from the larger remaining fold). For a candidate $(c,\theta)$ on grid
$\mathcal C\times\Theta$, the tune-fold objective approximates the expected region measure,
\begin{equation}
\bar W(c,\theta) = \frac{1}{|\mathcal D_{\mathrm{tune}}|}\sum_{i\in\mathcal D_{\mathrm{tune}}}
\big|\hat C_{c,\theta}(\bm x_i)\big|,
\end{equation}
% Why this sentence is spelled out explicitly: it is the concrete
% instantiation of the "tuned method can't be worse than its ancestor"
% claim from the Contributions paragraph -- because (1,0) (plain HPD) is
% literally a member of the search grid, the arg-min over that grid is, by
% definition, at least as good as evaluating the objective at (1,0) alone.
% This is what makes Proposition 1 more than a curiosity: it is the reason
% the tuning search has a built-in floor.
and $(c^\star,\theta^\star) = \arg\min_{(c,\theta)\in\mathcal C\times\Theta}\bar W(c,\theta)$.
Because $(1,0)\in\mathcal C\times\Theta$, $\bar W(c^\star,\theta^\star)\le \bar W(1,0)$ always holds
on the tune fold: STA-HPD's own tuning objective can never underperform plain HPD calibration
\emph{there}. Since a single split makes $(c^\star,\theta^\star)$ noticeably split-dependent, $n_{\mathrm{repeats}}$
(default $10$) independent tune/calib splits are drawn; $(c^\star,\theta^\star)$ is taken as the
per-split arg-min \emph{mode} (majority vote, appropriate for a discrete grid) rather than a mean,
and the final threshold is the \emph{median} of the per-split calib-fold quantiles at the
now-fixed $(c^\star,\theta^\star)$:
\begin{equation}
\hat t = \operatorname*{median}_{r=1,\dots,n_{\mathrm{repeats}}}\ \operatorname{Quantile}_{q_{\mathrm{level}}}
\Big(\big\{S_i(c^\star,\theta^\star)\big\}_{i\in\mathcal D_{\mathrm{calib}}^{(r)}}\Big).
\label{eq:medianquantile}
\end{equation}

% Why Proposition 2 is stated separately from Theorem 1/2: those two are
% about *width optimality*; this one is about *coverage validity*, which is
% a completely independent property in conformal prediction (a method can be
% valid but wide, or narrow but invalid). Keeping them as separate numbered
% results makes clear that STA-HPD's coverage guarantee does not depend on
% any of the optimality arguments holding -- it depends only on the tune/
% calib split being label-independent.
\begin{proposition}[Marginal validity]
\label{prop:validity}
For any $(c,\theta)$ fixed independently of the test label — in particular, one selected using
only $\mathcal D_{\mathrm{tune}}$ splits disjoint from the calib fold that supplies $\hat t$ —
Eq.~\eqref{eq:score} is a data-independent scoring function of $(\bm x_i,y_i)$, so the scores
$\{S_i(c,\theta)\}_{i\in\mathcal D_{\mathrm{calib}}}$ together with the test score are exchangeable,
and the quantile construction of Eq.~\eqref{eq:qlevel} yields
$P\big(y_{\mathrm{test}}\in \hat C(\bm x_{\mathrm{test}})\big)\ge 1-\alpha$ exactly, independent of
model or mixture correctness.
\end{proposition}

The proof is the standard split-conformal exchangeability argument \cite{vovk2005algorithmic},
applied conditionally on the (label-independent) tune split that fixes $(c^\star,\theta^\star)$.
Algorithm~\ref{alg:stahpd} summarizes the full fitting procedure.

% Why an explicit algorithm block, given the procedure is already described
% in prose above: the prose describes the *reasoning* for each step; the
% algorithm gives the exact *order of operations* (in particular, that the
% mode vote over (c,theta) happens across all repeats before the final
% median-quantile pass, not interleaved with it), which is the detail most
% likely to be implemented incorrectly from prose alone.
\begin{algorithm}[t]
\caption{STA-HPD fitting}
\label{alg:stahpd}
\begin{algorithmic}[1]
\State \textbf{Input:} calibration set $\mathcal D_{\mathrm{cal}}$, grids $\mathcal C,\Theta$, $\alpha$, $n_{\mathrm{repeats}}$
\For{$r = 1$ to $n_{\mathrm{repeats}}$}
  \State split $\mathcal D_{\mathrm{cal}}$ into $\mathcal D_{\mathrm{tune}}^{(r)}, \mathcal D_{\mathrm{calib}}^{(r)}$
  \For{$(c,\theta) \in \mathcal C\times\Theta$}
    \State compute $\bar W^{(r)}(c,\theta)$ on $\mathcal D_{\mathrm{tune}}^{(r)}$ (fixed-grid measure)
  \EndFor
  \State $(c^{(r)},\theta^{(r)}) \gets \arg\min \bar W^{(r)}$
\EndFor
\State $(c^\star,\theta^\star) \gets$ majority vote over $\{(c^{(r)},\theta^{(r)})\}_{r=1}^{n_{\mathrm{repeats}}}$
\For{$r = 1$ to $n_{\mathrm{repeats}}$}
  \State $\hat t^{(r)} \gets \operatorname{Quantile}_{q_{\mathrm{level}}}\big(\{S_i(c^\star,\theta^\star)\}_{i\in\mathcal D_{\mathrm{calib}}^{(r)}}\big)$
\EndFor
\State $\hat t \gets \operatorname{median}_r \hat t^{(r)}$
\State \textbf{Output:} $(c^\star,\theta^\star,\hat t)$; predict via Eq.~\eqref{eq:region}, roots located by grid scan $+$ Brent refinement
\end{algorithmic}
\end{algorithm}

% Why this closing paragraph of Section II describes the numerical root-
% finding machinery: eq:region defines the prediction set implicitly (as a
% level set), and unlike the earlier equations this is not something a
% reader can evaluate by inspection -- stating that it reduces to the same
% grid-scan-plus-Brent procedure as plain HPD (rather than something new)
% is what justifies the shared "predict cost" row later in Table I.
At prediction time, the exact set of Eq.~\eqref{eq:region} is located per test point by a coarse
grid scan for sign changes of $\log f_c(\cdot\mid\bm x) + g(\bm x)$ over a window built from all $K$
component tails, followed by bracketed root refinement (Brent's method) at each detected crossing;
adjacent segments separated by less than a numerical-noise tolerance are merged. This is identical
machinery to plain HPD calibration, applied to $f_{c^\star}$ instead of $f_{\mathrm{mix}}$, so
STA-HPD inherits plain HPD's degenerate reduction to a single symmetric interval whenever $K=1$ or
all components share a common mean, and its capacity to return disjoint multi-modal regions
otherwise.

\section{Comparative Analysis}
\label{sec:comparison}
% Why these three particular baselines (Standard CP, CP-VS, MoG-HPD) are
% named as "the" comparison set for the qualitative table below, even though
% Section IV's empirical table has more methods (MoECP, CQR, adaptive
% variance-ratio too): these three are the ones whose *functional form* is
% directly nested inside STA-HPD's family (per Proposition 1), so they are
% the right set for a theoretical (not just empirical) comparison table.
We compare STA-HPD against three standard baselines that share the split-conformal quantile step of
Eq.~\eqref{eq:qlevel} but differ in what they conformalize: Standard CP (raw absolute residual,
constant width), CP-VS (residual normalized by a single reported scale, symmetric heteroscedastic
width), and plain HPD/MoG-HPD (fixed-density level set, Theorem~\ref{thm:hpd}). Table~\ref{tab:comparison}
summarizes the comparison; CQR is included as a fourth, non-nested reference point since it is the
one baseline whose region shape is not a special case of Eq.~\eqref{eq:region}.

% Why table*[t] (full page width) rather than a single-column table: the
% five columns (method, fit cost, predict cost, free parameters, region
% shape) do not fit legibly in a single ~3.3in column at 10pt, so the table
% spans both columns; [t] places it at the top of the page it lands on,
% which is the usual convention for full-width floats in twocolumn layouts.
\begin{table*}[t]
\centering
\caption{Theoretical comparison of STA-HPD against baseline conformal calibrators. $n$: calibration
size, $K$: number of mixture components, $|\mathcal C|,|\Theta|$: grid sizes, $R$: number of tuning
repeats, $G_{\mathrm{grid}}$: per-point root-search grid resolution.}
\label{tab:comparison}
\begin{tabular}{@{}lllll@{}}
\toprule
\textbf{Method} & \textbf{Fit cost} & \textbf{Predict cost / point} & \textbf{Free parameters} & \textbf{Region shape} \\
\midrule
Standard CP & $O(n\log n)$ & $O(1)$ & $0$ & symmetric, constant width \\
CP-VS (total/aleatoric) & $O(n\log n)$ & $O(1)$ & $0$ & symmetric, $\propto \sigma(\bm x)$ \\
CQR & $O(n\log n)$ & $O(1)$ & $0$ (given quantile head) & asymmetric, additive shift \\
MoG-HPD & $O(n\log n) + O(K\,G_{\mathrm{grid}})$ & $O(K\,G_{\mathrm{grid}})$ & $0$ (density fixed) & possibly $K$ disjoint intervals \\
STA-HPD & $O\!\big(R\,|\mathcal C||\Theta|\,n_{\mathrm{tune}} + R\,n\log n\big)$ & $O(K\,G_{\mathrm{grid}})$ & $2$ ($c,\theta$), grid-selected & possibly $K$ disjoint intervals \\
\bottomrule
\end{tabular}
\end{table*}

% Why "Computational complexity" is its own labeled paragraph rather than
% folded into the table caption: the table gives the final Big-O only; the
% paragraph explains *why* STA-HPD's predict-time cost is identical to
% MoG-HPD's (same root-finding machinery, per the closing paragraph of
% Section II) even though its fit-time cost is higher -- a distinction that
% matters in practice (predict-time cost is paid once per deployment
% inference call; fit-time cost is paid once per calibration run).
\paragraph{Computational complexity.} Standard CP and CP-VS require only sorting the calibration
scores once. MoG-HPD adds, at both fit and predict time, a per-point root-finding pass over a grid
of resolution $G_{\mathrm{grid}}$ across $K$ component windows. STA-HPD's \emph{prediction} cost is
identical to MoG-HPD's, since the final region is located with the same grid-scan-plus-Brent
machinery applied to a single fixed $f_{c^\star}$; the additional cost is confined to \emph{fitting},
where the $(c,\theta)$ grid is scored on a subsampled tune fold ($n_{\mathrm{tune}} \le 500$ by
default) using a cheap vectorized fixed-grid measure rather than exact root-finding, since only the
arg-min over a discrete grid is needed there, not exact widths.

% Why "Parameter efficiency" needs its own paragraph: introducing two new
% tunable knobs could be read as adding overfitting risk; this paragraph
% pre-empts that objection by pointing back to Proposition 1 -- the knobs
% strictly enlarge what is reachable (they contain every baseline as a
% specific grid point) rather than replacing a zero-parameter method with an
% unconstrained one.
\paragraph{Parameter efficiency.} Standard CP, CP-VS, CQR, and MoG-HPD are all zero-parameter
calibrators in the sense that once the base model is frozen, the calibration step fits a single
scalar threshold with no further tunable structure. STA-HPD introduces exactly two scalar knobs,
each of which provably contains its ancestor method's behavior as a boundary value
(Proposition~\ref{prop:recovery}), so the extra parameters strictly enlarge the achievable family
rather than replacing it.

% Why "Convergence" is reinterpreted as "selection variance" rather than a
% classical optimization rate: (c*, theta*) come from a discrete grid
% arg-min, not gradient descent, so there is no gradient-based convergence
% rate to report -- this paragraph says so explicitly rather than silently
% omitting the requested "convergence guarantees" section item, and instead
% answers the question that actually matters for a grid search: how stable
% is the selected cell across resamples.
\paragraph{Convergence / stability of the tuned parameters.} Because $(c^\star,\theta^\star)$ are
selected from a finite discrete grid rather than by continuous gradient-based optimization, there is
no convergence rate in the classical optimization sense; the relevant stability question is
selection variance across resamples of the tune fold. A single $50/50$ tune/calib split was found
empirically to lose to plain HPD outright, since it starves $\hat t$ of data faster than the tuning
gain compensates. Repeated splitting with majority-vote aggregation over $R=10$ resamples
(Algorithm~\ref{alg:stahpd}) trades a small amount of fit cost for materially lower selection
variance, mirroring the cross-conformal-style stabilization used elsewhere for single-parameter
tuned conformal methods \cite{gibbs2021adaptive,tibshirani2019conformal}.

% Why this last comparison paragraph is included: the table only reports
% quantitative cost figures; the qualitative "why do these methods differ in
% region shape at all" question (symmetric interval vs. quantile-head shift
% vs. full density level set) is what actually explains *why* MoG-HPD and
% STA-HPD can win on multimodal data while Standard CP/CP-VS/CQR cannot,
% tying this section back to the motivation from the Introduction.
\paragraph{Architectural differences.} Standard CP and CP-VS commit to the symmetric-interval
functional form $\hat y(\bm x)\pm s(\bm x)$ irrespective of the model's density shape. CQR
delegates shape entirely to an auxiliary quantile-regression head and adds only a global additive
correction. MoG-HPD and STA-HPD both use the model's full mixture density and can return disjoint,
multi-modal regions; STA-HPD differs from MoG-HPD only in reshaping and rethresholding that density
\emph{before} conformalizing it, rather than accepting it as given.

\section{Experimental Simulation \& Results}
\label{sec:experiments}
% Why four specific datasets, and why the numbers below are what they are:
% these are not simulated for this paper -- they are the real, most recent
% calibration results already logged in this repository's
% result_calibration_*.txt files (same trained MOG checkpoints, same
% seed=4022, K=3 experts, run through each calibrator in turn), pulled
% directly rather than re-simulated, so every number in Tables II and III
% below is independently reproducible from those log files.
We evaluate on four tabular regression benchmarks with a $K=3$-expert MoE trained by mixture
negative log-likelihood (probabilistic-head, softmax-gated configuration, ``MOG''): a Synthetic
task whose true conditional law is a known trimodal mixture
$Y\mid X \sim 0.2\,\mathcal N(X,1)+0.3\,\mathcal N(X^2,1)+0.5\,\mathcal N(X^3,1)$; Bike
(bike-sharing demand); Temperature (bias-corrected temperature forecasting); and
Superconductivity (critical-temperature prediction). Nominal coverage is fixed at
$1-\alpha=90\%$; STA-HPD's calibration set is repeatedly split with $\texttt{tune\_frac}=0.2$ and
$R=10$ over the grid $\mathcal C=\{0.2,0.4,\dots,2.0\}$, $\Theta=\{-1.0,-0.75,\dots,3.0\}$.

% Why this paragraph states the coverage range before the table: it
% pre-empts the natural objection to any width comparison -- that a method
% could look narrower simply by under-covering. Reporting the achieved
% coverage band up front (88.6-91.5%) establishes that the width numbers in
% Table II are being compared at matched, not cherry-picked, coverage.
\paragraph{Full comparison against all baselines.} Table~\ref{tab:allmethods} reports average
interval width for every method introduced in Sections~\ref{sec:comparison}--\ref{sec:math} —
Standard CP, both variance-scaled CP variants, MoECP, CQR, and the width-tuned adaptive
variance-ratio calibrator, alongside MoG-HPD and STA-HPD — on all four datasets, at matched
coverage. Realized coverage across all $32$ (method, dataset) cells lies in $[88.6\%,91.5\%]$,
tightly centered on the $90\%$ nominal target, so the width comparison below is not confounded by
one method quietly under-covering to buy a narrower interval.

% Why \small is applied to this table specifically (unlike Table I): eight
% method rows by four dataset columns of 4-decimal numbers is denser than
% Table I's five-column qualitative table, and \small keeps it from
% overflowing the allotted table* width or forcing an awkward line break.
\begin{table*}[t]
\centering
\small
\caption{Average interval width by calibration method across all four datasets, $K{=}3$-expert MOG
architecture, nominal coverage $1-\alpha=90\%$. Narrowest width per dataset in \textbf{bold}.}
\label{tab:allmethods}
\begin{tabular}{@{}lrrrr@{}}
\toprule
\textbf{Method} & \textbf{Synthetic} & \textbf{Bike} & \textbf{Temperature} & \textbf{Superconductivity} \\
\midrule
Standard CP              & 2.4552 & 0.9465 & 1.0294 & 1.1740 \\
CP-VS (total)             & 1.3716 & 0.7926 & 1.0033 & 0.9475 \\
CP-VS (aleatoric)         & 1.6242 & \textbf{0.7713} & 0.9808 & 0.8588 \\
MoECP                     & 1.4914 & 0.9874 & 1.0287 & 0.9906 \\
CQR                       & 1.3570 & 1.0370 & 1.0809 & 0.9183 \\
Adaptive variance-ratio   & 1.3550 & 0.7720 & 0.9853 & 0.8647 \\
MoG-HPD                   & 1.0490 & 0.7940 & 0.9817 & 0.7964 \\
\textbf{STA-HPD}          & \textbf{0.9925} & 0.7729 & \textbf{0.9673} & \textbf{0.7816} \\
\bottomrule
\end{tabular}
\end{table*}

% Why this paragraph explicitly names Bike as an exception rather than only
% reporting the three wins: overstating "STA-HPD wins everywhere" would be
% contradicted by the reader's own inspection of Table II (CP-VS aleatoric
% and adaptive variance-ratio are both numerically narrower on Bike), so the
% honest framing is both more credible and consistent with how this
% repository's own calibration_methods.tex reports results elsewhere
% (it explicitly separates "real" gains from headline numbers). The
% Theorem 2 pointer here is what turns "Bike is an exception" from a bare
% observation into a predicted consequence of the theory.
STA-HPD is the narrowest method overall on three of the four datasets (Synthetic, Temperature,
Superconductivity), including a $5.4\%$ reduction over plain MoG-HPD on Synthetic — the dataset
whose true conditional law is genuinely multimodal — and a $1.9\%$ reduction on
Superconductivity, where it is also the narrowest interval among \emph{all} eight methods by a
comfortable margin over the next-best, CP-VS (aleatoric) at $0.8588$ ($-9.0\%$). Bike is the one
honest exception: STA-HPD ($0.7729$) is narrower than plain MoG-HPD ($0.7940$, $-2.7\%$) and every
other baseline except two — CP-VS (aleatoric) at $0.7713$ and the adaptive variance-ratio
calibrator at $0.7720$ — both of which it trails by under $0.2\%$, a gap within run-to-run
selection noise rather than a systematic loss. This is consistent with
Theorem~\ref{thm:optimality}: Bike is the dataset where a symmetric, purely aleatoric-scaled
interval is already close to the population-optimal shape, leaving little for the density-level-set
family to add. Standard CP and, on three of four datasets, MoECP are the widest methods throughout,
confirming that ignoring the model's own uncertainty structure (Standard CP) or relying on routing
similarity alone without a variance or density model (MoECP) leaves the most width on the table.

% Why a second, smaller table repeats two columns already in Table II:
% Table II answers "how does STA-HPD compare to everything"; this one
% answers a narrower, theory-specific question -- "does the one comparison
% Proposition 1 makes a guarantee about actually hold" -- and adds the
% (c*, theta*) column that would clutter the eight-row Table II. Splitting
% them keeps each table answering one question.
\paragraph{STA-HPD vs.\ its immediate ancestor.} Table~\ref{tab:results} isolates the comparison
that Proposition~\ref{prop:recovery} makes a guarantee about: STA-HPD against the plain HPD
calibrator it strictly contains at $(c,\theta)=(1,0)$.

\begin{table}[H]
\centering
\caption{MoG-HPD vs.\ STA-HPD at matched coverage, with the grid-selected $(c^\star,\theta^\star)$.}
\label{tab:results}
\begin{tabular}{@{}lrrrl@{}}
\toprule
\textbf{Dataset} & \textbf{MoG-HPD} & \textbf{STA-HPD} & $\Delta$ & $(c^\star,\theta^\star)$ \\
\midrule
Synthetic & 1.0490 & \textbf{0.9925} & $-5.4\%$ & $(1.3, 0.25)$ \\
Bike & 0.7940 & \textbf{0.7729} & $-2.7\%$ & $(1.0, 0.75)$ \\
Temperature & 0.9817 & \textbf{0.9673} & $-1.5\%$ & $(1.3, 0.25)$ \\
Superconductivity & 0.7964 & \textbf{0.7816} & $-1.9\%$ & $(0.8, 0.25)$ \\
\bottomrule
\end{tabular}
\end{table}

% Why the fitted c* values are discussed individually rather than just
% reported: c*>1 on two datasets is, at first glance, the opposite of the
% "shrinkage" intuition argued in Section II-D -- this paragraph resolves
% that apparent tension by pointing out Theorem 2 only predicts the
% *direction* of correction when sigma_k is under-informative, and the
% search is free to go either way depending on the dataset, which is
% exactly what a two-sided knob (rather than a one-sided shrinkage-only
% method) is supposed to allow.
The reduction is positive on all four datasets, as Proposition~\ref{prop:recovery} guarantees on
the tune fold, and never reverses on the held-out calibration/test split either. The fitted shape
exponent $c^\star>1$ on Synthetic and Temperature indicates the search found it worthwhile to
\emph{amplify}, not shrink, the model's reported heteroscedasticity at those two datasets, while
$c^\star=0.8<1$ on Superconductivity matches the shrinkage direction predicted by
Theorem~\ref{thm:optimality} when $\sigma_k(\bm x)$ is only partially informative about the
residual. No dataset reduces to $(c^\star,\theta^\star)=(1,0)$, i.e.\ the search never collapses
back onto plain HPD calibration, consistent with Proposition~\ref{prop:recovery}'s guarantee that
the tuning objective can only improve on or match it.

% Why this new section is inserted here, as a self-contained addition, rather than reworking
% Sections II-IV in place: STA-HPD's own results above are unchanged and already published-quality;
% the cleanest way to add a genuinely new degree of freedom is to state exactly which part of the
% MoG variance decomposition it was still missing, extend the same surrogate-family/validity/
% boundary-recovery machinery by one exponent, and report its own matched comparison against the
% method it extends -- STA-HPD itself, not just plain HPD. This mirrors how Section II already
% introduced STA-HPD as a strict extension of MoG-HPD.
\section{SETA-HPD: Reshaping the Between-Component (Epistemic) Variance}
\label{sec:seta}

\subsection{The gap STA-HPD leaves open}
% Why lead with the variance-decomposition identity again: it is the organizing fact for this
% whole section -- STA-HPD's own sigma_tot (Sec. II-D) already writes out both terms, but only ever
% reshapes the first one (via tilde_sigma_k). Naming that omission explicitly is what motivates the
% new knob, exactly as the Introduction's two failure modes motivated c and theta in the first place.
Every predictive law in this family already decomposes, by the law of total variance, as
\begin{equation}
\operatorname{Var}(Y\mid\bm x) = \underbrace{\textstyle\sum_k w_k\sigma_k^2}_{\text{within-component (aleatoric)}}
\;+\; \underbrace{\textstyle\sum_k w_k(\mu_k-\hat y)^2}_{\text{between-component (epistemic)}},
\label{eq:lotv}
\end{equation}
and $\sigma_{\mathrm{tot}}^2(\bm x;c)$ of Sec.~\ref{sec:math} already contains both terms. STA-HPD's
shape exponent $c$, however, only ever reshapes the first (via $\tilde\sigma_k(\bm x;c)$); the second
term -- how far apart the gate places the expert means -- is left exactly as the base model reported
it, however systematically over- or under-separated that reported spread is. This matters in
practice: the load-balancing penalty this repo adds to every MoE loss (Sec.~\ref{sec:experiments};
see also Exp\_Tabular\_Regression.\_calculate\_loss) is needed to keep the gate from collapsing onto
a single expert, but the same pressure that prevents collapse can push well-separated experts further
apart than the true conditional law warrants, opening exactly the kind of artificially deep density
valley that Example~\ref{ex:bimodal} shows is costly for a level-set method to cover. SETA-HPD
(Shape--Epistemic--Threshold Adaptive HPD) adds the missing exponent.

\subsection{The SETA-HPD family}
% Why the mean-shift is introduced the same way Sec. II-D introduced tilde_sigma_k (define the
% transform, then a lemma stating exactly what it preserves and what it scales): symmetry with
% Lemma 1 is deliberate -- a reader who accepted Lemma 1's affine-in-log-scale argument for c should
% be able to check this construction's two claims (mean preservation, variance scaling) just as
% mechanically, without re-deriving a new proof style from scratch.
Introduce a second exponent $\rho\ge0$ acting on the between-component term only, via the mean-shift
\begin{equation}
\tilde\mu_k(\bm x;\rho) = \hat y(\bm x) + \sqrt{\rho}\,\big(\mu_k(\bm x)-\hat y(\bm x)\big),
\label{eq:etamu}
\end{equation}
and define the jointly-reshaped surrogate density and total scale
\begin{align}
f_{c,\rho}(y\mid\bm x) &= \sum_{k=1}^K w_k(\bm x)\,\mathcal N\big(y;\ \tilde\mu_k(\bm x;\rho),\ \tilde\sigma_k^2(\bm x;c)\big), \label{eq:fcrho}\\
\sigma_{\mathrm{tot}}^2(\bm x;c,\rho) &= \sum_k w_k\,\tilde\sigma_k^2(\bm x;c) \;+\; \rho\sum_k w_k\,(\mu_k-\hat y)^2. \label{eq:sigmatotrho}
\end{align}

\begin{lemma}[The mean-shift preserves $\hat y$ and rescales the epistemic term exactly]
\label{lem:etashift}
For every $\bm x$, $\sum_k w_k\tilde\mu_k(\bm x;\rho) = \hat y(\bm x)$, and
$\sum_k w_k\big(\tilde\mu_k(\bm x;\rho)-\hat y(\bm x)\big)^2 = \rho\sum_k w_k(\mu_k(\bm x)-\hat y(\bm x))^2$.
Consequently $\tilde\mu_k(\bm x;1)=\mu_k(\bm x)$ (the identity) and $\tilde\mu_k(\bm x;0)=\hat y(\bm x)$
for every $k$ (every expert collapses onto the aggregate prediction).
\end{lemma}
\begin{proof}
Both identities follow by substituting Eq.~\eqref{eq:etamu} directly:
$\sum_k w_k\tilde\mu_k = \sum_k w_k\hat y + \sqrt\rho\sum_k w_k(\mu_k-\hat y) = \hat y + \sqrt\rho\,(\hat y - \hat y) = \hat y$,
using $\sum_k w_k=1$ and the definition $\hat y=\sum_k w_k\mu_k$; and
$\sum_k w_k(\tilde\mu_k-\hat y)^2 = \sum_k w_k\big(\sqrt\rho(\mu_k-\hat y)\big)^2 = \rho\sum_k w_k(\mu_k-\hat y)^2$
directly. Substituting $\rho=1,0$ gives the two boundary identities.
\end{proof}

% Why this remark, immediately after the lemma: it is the one-sentence payoff that justifies
% calling rho a knob on the epistemic term *specifically* -- without it, a reader would have to
% infer from Eqs. (24)-(25) alone that c and rho act on disjoint pieces of Eq. (23), rather than
% having it stated as a direct consequence of Lemma 2.
\begin{remark}
By Lemma~\ref{lem:etashift} and the definition of $\tilde\sigma_k(\bm x;c)$ (Lemma~\ref{lem:logaffine}),
$c$ and $\rho$ act on disjoint terms of the decomposition Eq.~\eqref{eq:lotv}: $c$ reshapes only the
aleatoric term, $\rho$ reshapes only the epistemic term, and neither changes the aggregate
prediction $\hat y(\bm x)$. $\rho<1$ compresses the experts toward $\hat y$ (filling exactly the kind
of density valley described above); $\rho>1$ amplifies their separation; $\rho=1$ recovers STA-HPD's
own $f_c$ and $\sigma_{\mathrm{tot}}(\bm x;c)$ exactly.
\end{remark}

The score, threshold, and induced region are the direct three-exponent analogues of
Eqs.~\eqref{eq:score}--\eqref{eq:region},
\begin{align}
S_i(c,\rho,\theta) &= -\log f_{c,\rho}(y_i\mid\bm x_i) - \theta\log\sigma_{\mathrm{tot}}(\bm x_i;c,\rho), \label{eq:setascore}\\
\hat C(\bm x) &= \Big\{y:\log f_{c,\rho}(y\mid\bm x)\ge -\hat t(c,\rho,\theta) - \theta\log\sigma_{\mathrm{tot}}(\bm x;c,\rho)\Big\}. \label{eq:setaregion}
\end{align}

\begin{proposition}[Boundary recovery, three exponents]
\label{prop:setarecovery}
$(c,\rho,\theta)=(1,1,0)$ recovers plain HPD calibration exactly. $(c,\rho,\theta)=(c,1,\theta)$ for
any $(c,\theta)$ recovers STA-HPD exactly, so STA-HPD is a full sub-family of SETA-HPD, not merely a
single point. $(c,\rho,\theta)=(1,\rho,0)$ for any $\rho$ recovers an epistemic-only HPD calibrator
that scales the between-component spread without reshaping the aleatoric term or tilting the
threshold. $(c,\rho,\theta)=(1,1,1)$ recovers CP-VS on the total predictive standard deviation
exactly when $K=1$ (where the epistemic term vanishes identically, so $\rho$ has no effect).
\end{proposition}
\begin{proof}
Each case is direct substitution into Eqs.~\eqref{eq:etamu}--\eqref{eq:setaregion} using
Lemma~\ref{lem:etashift} ($\tilde\mu_k(\bm x;1)=\mu_k$) and Lemma~\ref{lem:logaffine}
($\tilde\sigma_k(\bm x;1)=\sigma_k$), identically to the argument for
Proposition~\ref{prop:recovery}.
\end{proof}

Because $(1,0)\in\mathcal C\times\Theta$ is on STA-HPD's own grid and $\rho=1$ is always kept on
SETA-HPD's $\rho$-grid, the tuning objective's minimum over $\mathcal C\times\mathcal R\times\Theta$
can never exceed the minimum over the $\rho=1$ slice, which is STA-HPD's own tuning objective: SETA-HPD
cannot be worse than STA-HPD -- or, by the same argument through $(1,1,0)$, than plain HPD -- on the
tune fold. Validity is unaffected: Proposition~\ref{prop:validity}'s argument applies verbatim with
$(c,\rho,\theta)$ in place of $(c,\theta)$, since nothing about the exchangeability argument depends
on how many exponents the frozen, label-independent scoring function has.

\subsection{Fitting}
% Why this subsection stays short (no new Algorithm float): Algorithm 1 already describes the
% majority-vote/median-quantile procedure in full; the only structural change here is the grid's
% dimensionality, which is a one-line description, not a new control-flow diagram. Re-drawing the
% whole algorithm for one added inner loop would repeat Algorithm 1 without adding information.
Fitting extends Algorithm~\ref{alg:stahpd} by widening the grid searched at each split from
$\mathcal C\times\Theta$ to $\mathcal C\times\mathcal R\times\Theta$, with $\mathcal R=\{0,0.25,0.5,
0.75,1.0,1.5,2.0\}$ by default (1.0 always included). Since $f_{c,\rho}$ does not depend on $\theta$,
the tune-fold objective $\bar W(c,\rho,\theta)$ is evaluated by computing the reshaped density once
per $(c,\rho)$ pair and sweeping $\theta$ against it with only a quantile and a threshold-count per
value, rather than re-evaluating the mixture density for every $(c,\rho,\theta)$ triple; this keeps
the added cost of the third exponent to roughly a constant factor in $|\mathcal R|$ rather than a
factor in $|\mathcal R|\cdot|\Theta|$.

\subsection{Results}
% Why this table compares against STA-HPD specifically, not (only) plain MoG-HPD: Table III already
% answers "does STA-HPD beat MoG-HPD"; this table answers the question this section actually adds --
% "does reshaping the epistemic term on top of STA-HPD's own tuned (c,theta) buy anything further",
% which is a strictly harder bar to clear than beating the untuned baseline. rho* is reported because
% Proposition 3's boundary-recovery claim is only informative if the reader can see whether the
% search actually leaves rho=1.
Table~\ref{tab:seta} reports SETA-HPD against both plain MoG-HPD and STA-HPD, same checkpoints,
$1-\alpha=90\%$, single tune/calib split ($\texttt{split\_seed}=0$) for both STA-HPD and SETA-HPD so
the two are compared on identical folds.

\begin{table}[H]
\centering
\caption{MoG-HPD vs.\ STA-HPD vs.\ SETA-HPD at matched coverage. $\rho^\star$ is SETA-HPD's selected
epistemic exponent; $\rho^\star=1$ means the search left the epistemic term untouched.}
\label{tab:seta}
\begin{tabular}{@{}lrrrrr@{}}
\toprule
\textbf{Dataset} & \textbf{MoG-HPD} & \textbf{STA-HPD} & \textbf{SETA-HPD} & $\Delta$ & $\rho^\star$ \\
\midrule
Synthetic & 1.0490 & 0.9925 & 0.9925 & $0.0\%$ & $1.00$ \\
Bike & 0.7940 & 0.7729 & \textbf{0.7399} & $-4.3\%$ & $0.25$ \\
Temperature & 0.9817 & 0.9673 & 0.9673 & $0.0\%$ & $1.00$ \\
Superconductivity & 0.7964 & 0.7816 & \textbf{0.7849} & $+0.4\%$ & $0.75$ \\
\bottomrule
\end{tabular}
\end{table}

% Why the discussion is this candid rather than framed as a uniform win: the honest empirical
% picture (verified against 5 independent tune/calib splits per dataset, not just the single split
% in Table IV) is that rho* is stable but its practical payoff is dataset-dependent and, on two of
% four datasets, exactly zero -- reporting that plainly is what Proposition 3's "cannot be worse"
% guarantee is actually for: it is a floor, not a promise that the extra parameter always helps.
$\Delta$ is never positive relative to STA-HPD by more than measurement noise, consistent with
Proposition~\ref{prop:setarecovery}'s guarantee, but the practical picture is mixed rather than a
uniform win. On \textbf{Synthetic} and \textbf{Temperature}, $\rho^\star=1$ is selected and stable
across 5 independent tune/calib splits: STA-HPD's own $(c,\theta)$ already capture the correctable
mismatch on these two datasets, and there is no further epistemic miscalibration left for $\rho$ to
fix. On \textbf{Bike}, $\rho^\star=0.25$ is selected on all 5 splits (fully stable) and yields the
largest gain in the sweep, $-4.3\%$ width relative to STA-HPD, while eliminating disjoint
multi-segment sets entirely (from $1.4\%$ of test points to $0\%$): compressing the two experts
toward $\hat y$ removes the density valley that was fragmenting some predictions into two intervals.
On \textbf{Superconductivity}, $\rho^\star=0.75$ is likewise selected on all 5 splits; mean width is
flat, but the fraction of fragmented sets drops from $6.9\%$ to $4.2\%$ -- the same coverage budget
is spent as fewer, more usable single intervals rather than as raw width reduction. In no case does
SETA-HPD's test-set width exceed STA-HPD's by more than seed-to-seed noise, matching
Proposition~\ref{prop:setarecovery}'s tuning-fold guarantee; STA-HPD and SETA-HPD are implemented as
two separate, independently runnable calibrators (differing only in whether $\rho$ is searched), so
either can be deployed on its own.

% Why these two sections are appended here, in the same style as Sec. V: SETA-HPD already
% established the pattern this paper uses to extend the family -- name the exact gap left open by
% the previous member, extend the surrogate-family/validity/boundary-recovery machinery by exactly
% one new mechanism, and report a matched comparison against the method being extended, not just
% against plain HPD. EASE-HPD and MELD-HPD are two independent extensions of SETA-HPD (one widens
% the score by a fourth exponent, the other only changes how (c,rho) are SELECTED), so they get two
% short, focused sections rather than being folded into Sec. V, which is already a complete,
% self-contained unit.
\section{EASE-HPD: Reshaping the Threshold by Epistemic Share}
\label{sec:ease}

\subsection{The gap SETA-HPD leaves open}
% Why lead with sigma_tot's own decomposition: eq:sigmatotrho already splits into an aleatoric and
% an epistemic piece, but theta (both in STA-HPD and SETA-HPD) only ever sees their SUM through
% log sigma_tot -- naming that explicitly is what motivates a threshold field that looks at the
% MIX, not just the total, exactly as Sec. V-A named the missing rho by pointing at eq:lotv.
SETA-HPD's threshold field $\theta$ tilts the conformal threshold by the point's total predictive
scale $\sigma_{\mathrm{tot}}(\bm x;c,\rho)$ of Eq.~\eqref{eq:sigmatotrho}, which is itself a sum of an
aleatoric term and an epistemic term. Write these two terms explicitly,
\begin{equation}
A(\bm x;c) = \sum_{k=1}^K w_k(\bm x)\,\tilde\sigma_k^2(\bm x;c), \quad
E(\bm x;\rho) = \sum_{k=1}^K w_k(\bm x)\,\big(\tilde\mu_k(\bm x;\rho)-\hat y(\bm x)\big)^2,
\label{eq:aeterms}
\end{equation}
so that $\sigma_{\mathrm{tot}}^2(\bm x;c,\rho) = A(\bm x;c) + E(\bm x;\rho)$ restates
Eq.~\eqref{eq:sigmatotrho} (using Lemma~\ref{lem:etashift} to write the epistemic term as
$E(\bm x;\rho)=\rho\sum_k w_k(\mu_k(\bm x)-\hat y(\bm x))^2$). Because $\theta$ multiplies
$\log\sigma_{\mathrm{tot}}$, two points with the same total $A(\bm x;c)+E(\bm x;\rho)$ but opposite
\emph{compositions} -- one driven entirely by expert disagreement, the other entirely by
within-expert noise -- receive an identical threshold tilt, since $\theta$ only ever observes the
sum. This matters because the two terms are miscalibrated for structurally different reasons in
practice: the epistemic term is shaped by the load-balancing pressure named in
Sec.~\ref{sec:seta}-A, an artifact of gate regularization rather than a property of the outcome,
while the aleatoric term is fit directly by the per-expert negative log-likelihood and is
comparatively well behaved already. EASE-HPD (Epistemic-Aleatoric Share-adaptive threshold HPD) adds
a fourth exponent that lets the threshold field respond to the \emph{composition} of the variance,
not only its total.

\subsection{The epistemic share and the EASE-HPD family}
% Why the share is defined as a ratio in [0,1], and why a Lemma is needed before the score is
% written down: unlike c and rho, which reshape the DENSITY itself, phi tilts the threshold by a
% quantity -- s(x) -- that must be well-defined and bounded for every (c,rho) cell on the grid,
% including the boundary cells (rho=0, or K=1) where the epistemic term vanishes identically. The
% Lemma states exactly when that happens, which is what later explains, rather than merely asserts,
% why phi is a documented no-op in exactly those regimes.
Define the (surrogate, i.e.\ post-$c$, post-$\rho$) \emph{epistemic share}
\begin{equation}
s(\bm x;c,\rho) = \frac{E(\bm x;\rho)}{A(\bm x;c)+E(\bm x;\rho)} = \frac{E(\bm x;\rho)}{\sigma_{\mathrm{tot}}^2(\bm x;c,\rho)} \in [0,1].
\label{eq:share}
\end{equation}

\begin{lemma}[Epistemic share is well defined and collapses exactly when the epistemic term does]
\label{lem:share}
For every $\bm x,c,\rho$, $A(\bm x;c)\ge0$ and $E(\bm x;\rho)\ge0$, so $s(\bm x;c,\rho)\in[0,1]$ is
well defined whenever $\sigma_{\mathrm{tot}}(\bm x;c,\rho)>0$, which holds since every
$\tilde\sigma_k(\bm x;c)>0$. Moreover $s(\bm x;c,\rho)\equiv 0$ for every $\bm x$ whenever $\rho=0$
or $K=1$.
\end{lemma}
\begin{proof}
$A,E\ge0$ is immediate, since each is a weighted sum of squares with weights $w_k\ge0$. At $\rho=0$,
Lemma~\ref{lem:etashift} gives $\tilde\mu_k(\bm x;0)=\hat y(\bm x)$ for every $k$, so
$E(\bm x;0)=\sum_k w_k(\hat y(\bm x)-\hat y(\bm x))^2=0$ for every $\bm x$. At $K=1$,
$\hat y(\bm x)=\mu_1(\bm x)$ identically, so $E(\bm x;\rho)=w_1(\bm x)(\tilde\mu_1(\bm x;\rho)-\hat y(\bm x))^2=0$
for every $\bm x$ and every $\rho$, again by Lemma~\ref{lem:etashift} ($\tilde\mu_1$ is then a fixed
point of the shift, since there is nothing to shift toward). In both cases $s(\bm x;c,\rho)=0/A(\bm x;c)=0$.
\end{proof}

Freeze the calibration-set mean share at each $(c,\rho)$ cell, $\bar s(c,\rho) =
\tfrac1n\sum_{i=1}^n s(\bm x_i;c,\rho)$, exactly as $G$ is frozen for the shape exponent. A fourth
exponent $\phi\in\mathbb R$ then tilts the score by the point's own \emph{deviation} of its share from
that typical value,
\begin{align}
S_i(c,\rho,\theta,\phi) &= -\log f_{c,\rho}(y_i\mid\bm x_i) - \theta\log\sigma_{\mathrm{tot}}(\bm x_i;c,\rho) \notag\\
&\qquad{}- \phi\big(s(\bm x_i;c,\rho)-\bar s(c,\rho)\big), \label{eq:easescore}
\end{align}
with induced region
\begin{align}
\hat C(\bm x) = \Big\{y :\ \log f_{c,\rho}(y\mid\bm x) &\ge -\hat t(c,\rho,\theta,\phi) - \theta\log\sigma_{\mathrm{tot}}(\bm x;c,\rho) \notag\\
&\qquad{}- \phi\big(s(\bm x;c,\rho)-\bar s(c,\rho)\big)\Big\}. \label{eq:easeregion}
\end{align}
Centering by $\bar s(c,\rho)$ keeps $\phi$ and $\theta$ from fighting over the same additive constant
in $\hat t$, and boundedness of $s$ in $[0,1]$ (Lemma~\ref{lem:share}) rules out the log-ratio blow-ups
a raw ratio $E/A$ would risk when $A\to0$.

\begin{proposition}[Boundary recovery, four exponents]
\label{prop:easerecovery}
$(c,\rho,\theta,\phi)=(1,1,0,0)$ recovers plain HPD calibration exactly. $(c,\rho,\theta,0)$ for any
$(c,\rho,\theta)$ recovers SETA-HPD exactly, so SETA-HPD -- and, by
Proposition~\ref{prop:setarecovery}, STA-HPD and plain HPD calibration with it -- are sub-families
of EASE-HPD, not merely isolated grid points.
\end{proposition}
\begin{proof}
At $\phi=0$ the term $-\phi(s(\bm x;c,\rho)-\bar s(c,\rho))$ vanishes identically in
Eqs.~\eqref{eq:easescore}--\eqref{eq:easeregion}, regardless of the value of $s$, leaving exactly
SETA-HPD's score and region (Eqs.~\eqref{eq:setascore}--\eqref{eq:setaregion}); the case
$(1,1,0,0)$ then follows from Proposition~\ref{prop:setarecovery}'s own $(1,1,0)$ case.
\end{proof}

\begin{remark}[$\phi$ is provably inert exactly where Lemma~\ref{lem:share} predicts]
\label{rem:phiinert}
Whenever the tuned $\rho^\star=0$ or $K=1$, Lemma~\ref{lem:share} gives $s(\bm x;c,\rho)\equiv0$ for
every $\bm x$ in the calibration set, hence $\bar s(c,\rho)=0$ and $s(\bm x;c,\rho)-\bar s(c,\rho)\equiv0$
identically: the $\phi$ term contributes exactly zero to Eqs.~\eqref{eq:easescore}--\eqref{eq:easeregion}
for \emph{every} value of $\phi$, so $\phi$ is unidentifiable from the tune-fold objective and any tie
among $\phi$-values is broken arbitrarily. This is not a failure of the mechanism; it is the direct,
checkable consequence of Lemma~\ref{lem:share}, and is exactly what the worked example below
exhibits.
\end{remark}

Validity is unaffected: Proposition~\ref{prop:validity}'s argument applies verbatim with
$(c,\rho,\theta,\phi)$ in place of $(c,\theta)$, since $\phi(s(\bm x;c,\rho)-\bar s(c,\rho))$ is, for
any $(c,\rho,\theta,\phi)$ fixed on a tune split disjoint from the calib fold, a data-independent
function of $\bm x$ alone, so the resulting score remains a frozen, label-independent function of
$(\bm x_i,y_i)$ by the time it is conformalized.

\subsection{Fitting}
Fitting widens Algorithm~\ref{alg:stahpd}'s search from $\mathcal C\times\mathcal R\times\Theta$ to
$\mathcal C\times\mathcal R\times\Theta\times\Phi$, with $\Phi=\{-1.0,-0.5,0,0.5,1.0\}$ by default
($0\in\Phi$ always). Since $f_{c,\rho}$ and $s(\cdot;c,\rho)$ depend on neither $\theta$ nor $\phi$,
both the density and the share are computed once per $(c,\rho)$ pair and the whole $(\theta,\phi)$
sub-grid is swept with only quantile and threshold-count operations, mirroring the
$\theta$-factorization already used for SETA-HPD's $(c,\rho,\theta)$ grid. Because $\rho=1$ and
$\phi=0$ are always on their respective grids, $\bar W(c^\star,\rho^\star,\theta^\star,\phi^\star) \le
\bar W(c,\rho,\theta,0) \le \bar W(1,1,0,0)$ on the tune fold for any $(c,\rho,\theta)\in\mathcal
C\times\mathcal R\times\Theta$: EASE-HPD's tuning objective can never exceed SETA-HPD's, which can
never exceed STA-HPD's or plain HPD's -- the same nested floor argument as
Sec.~\ref{sec:seta}, extended by one more axis.

\subsection{A worked example}
% Why this subsection reports a single config rather than a multi-dataset table like Table III/IV:
% the result logs currently checked into this repository for EASE-HPD and MELD-HPD cover one
% checkpoint (the single-expert Synthetic MOG model, seed 4022); reporting only what those logs
% actually contain, rather than re-deriving or estimating numbers for datasets not present in them,
% keeps every figure in this subsection independently reproducible from
% result_calibration_ease_hpd.txt. What the config lacks in breadth it makes up for as a clean
% illustration of Remark~\ref{rem:phiinert}, since it is exactly the collapse case the remark
% describes.
On the single-expert Synthetic MOG checkpoint (seed $4022$, $1-\alpha=90\%$) already used to
illustrate STA-HPD's own boundary case in Table~\ref{tab:results}, the search selects
$(c^\star,\rho^\star,\theta^\star,\phi^\star)=(0.8,0,1.5,-1.0)$, with tune-fold width identical across
every $\phi\in\Phi$ (the logged $\phi$-marginal is $1.4402$ at all five grid values). This is exactly
Remark~\ref{rem:phiinert}'s collapse case: the search independently selects $\rho^\star=0$ (matching
SETA-HPD's own choice on this config, Table~\ref{tab:seta}'s pattern), which by
Lemma~\ref{lem:share} forces $s(\bm x;c,0)\equiv0$ and $\bar s=0$, so $\phi$ is provably unable to move
either the score or the region here regardless of its selected value. EASE-HPD's realized coverage
and width, $0.8933$ and $1.3473$, are accordingly identical to STA-HPD's and SETA-HPD's on this same
config (Table~\ref{tab:allfive} below) to four decimals -- an empirical confirmation, not merely a
theoretical prediction, of Remark~\ref{rem:phiinert}.

% Why this new section is titled around "proper scoring rule" rather than "log-likelihood
% selection": that phrasing is what lets the KL-identity Proposition below, and the honesty-forcing
% Remark that follows it, both attach to a named, citable statistical concept (Gneiting & Raftery's
% strictly-proper-scoring-rule framework) rather than reading as an ad hoc heuristic.
\section{MELD-HPD: Selecting Shape by a Proper Scoring Rule}
\label{sec:meld}

\subsection{Why tune-fold width can misfire as a selector}
% Why this subsection is candid about a real empirical failure mode before proposing the fix:
% Sec. V and the EASE-HPD section above both lean on the "grid contains its ancestor, so the tuning
% objective can only improve" floor -- but that floor is a fact about a NOISY, subsampled estimator
% of population width, not about population width itself, and docs/sta_hpd_mog_decomposition.md
% already documents SETA-HPD losing to STA-HPD out of sample despite the floor holding in-sample.
% Naming that explicitly is what motivates using a different, smoother objective for (c,rho).
STA-HPD, SETA-HPD, and EASE-HPD all select their density-shaping exponents by $\arg\min$ of the
tune-fold width estimator $\bar W$ (Sec.~\ref{sec:math}-E), a Monte-Carlo approximation of
$\mathbb E_{\bm x}[|C(\bm x)|]$ evaluated on a subsampled tune fold of at most $\mathtt{max\_tune\_points}$
points. Widening the grid searched by $\bar W$'s $\arg\min$ gives it more room to fit sampling noise
in that specific estimator, not necessarily to reduce the population width it approximates; the
$\rho$ selected this way does not always generalize; SETA-HPD is empirically documented to come out
\emph{wider} than STA-HPD on held-out test folds on some configurations despite
$\bar W(c^\star,\rho^\star,\theta^\star)\le\bar W(1,1,0)$ holding exactly on the tune fold that chose
it. The floor Sec.~\ref{sec:seta} and Sec.~\ref{sec:ease} establish is real, but it is a floor on the
\emph{estimator}, not on the population quantity the estimator targets. MELD-HPD
(Maximum-likelihood Epistemic-aLeatoric Decomposition) selects the two \emph{density-shaping}
exponents $(c,\rho)$ -- the only ones that change $f_{c,\rho}$ itself -- by a different, smooth, and
strictly proper objective instead, evaluated on the entire tune fold rather than a width subsample.

\subsection{The log-score objective and what it targets}
Define the mean calibration log-density on the tune fold,
\begin{equation}
L(c,\rho) = \frac{1}{|\mathcal D_{\mathrm{tune}}|}\sum_{i\in\mathcal D_{\mathrm{tune}}} \log f_{c,\rho}(y_i\mid\bm x_i).
\label{eq:logscore}
\end{equation}

\begin{proposition}[The log score targets the KL-closest member of the family]
\label{prop:logscore-kl}
Let $p(\cdot\mid\bm x)$ denote the true conditional density of $Y\mid X=\bm x$, and let
$\mathbb E$ denote expectation under the true law of $(X,Y)$. Then
\begin{equation}
\mathbb E\big[\log f_{c,\rho}(Y\mid X)\big] = -\,\mathbb E_X\big[H(p(\cdot\mid X))\big] \;-\; \mathbb E_X\big[\mathrm{KL}\big(p(\cdot\mid X)\,\big\|\,f_{c,\rho}(\cdot\mid X)\big)\big],
\label{eq:klidentity}
\end{equation}
where $H$ is differential entropy and $\mathrm{KL}$ is Kullback--Leibler divergence. The first term
does not depend on $(c,\rho)$, so
$\arg\max_{c,\rho}\ \mathbb E[\log f_{c,\rho}(Y\mid X)] = \arg\min_{c,\rho}\ \mathbb E_X[\mathrm{KL}(p(\cdot\mid X)\,\|\,f_{c,\rho}(\cdot\mid X))]$:
the population log score is maximized exactly at whichever $(c,\rho)$ makes the surrogate density
closest, on average in KL divergence, to the true conditional law.
\end{proposition}
\begin{proof}
For fixed $\bm x$, write $\int p(y\mid\bm x)\log f_{c,\rho}(y\mid\bm x)\,dy = \int p(y\mid\bm x)\log
p(y\mid\bm x)\,dy - \int p(y\mid\bm x)\log\big(p(y\mid\bm x)/f_{c,\rho}(y\mid\bm x)\big)\,dy = -H(p(\cdot\mid\bm x))
- \mathrm{KL}(p(\cdot\mid\bm x)\,\|\,f_{c,\rho}(\cdot\mid\bm x))$, the standard cross-entropy
decomposition. Taking $\mathbb E_X$ of both sides and using
$\mathbb E[\log f_{c,\rho}(Y\mid X)] = \mathbb E_X\big[\int p(y\mid X)\log f_{c,\rho}(y\mid X)\,dy\big]$
gives Eq.~\eqref{eq:klidentity}. Since $\mathrm{KL}\ge0$ with equality iff the two densities agree a.e.,
and $\mathbb E_X[H(p(\cdot\mid X))]$ is a fixed constant of the true law, maximizing the left side over
$(c,\rho)$ is exactly minimizing $\mathbb E_X[\mathrm{KL}(p(\cdot\mid X)\|f_{c,\rho}(\cdot\mid X))]$.
\end{proof}

This is the population-level statement of why the log score is a \emph{strictly proper} scoring
rule \cite{gneiting2007strictly}: no other density in the family can achieve a higher expected score
than the KL-closest one, so maximizing the sample analogue $L(c,\rho)$ of Eq.~\eqref{eq:logscore} is
a consistent estimator of the population-optimal shape as $|\mathcal D_{\mathrm{tune}}|\to\infty$,
using the full tune fold rather than a width subsample, and with no dependence on grid resolution
for the underlying quantity being estimated (unlike $\bar W$, which additionally discretizes the
region measure).

\begin{remark}[Log score is a proxy for, not a guarantee of, width optimality]
\label{rem:logscoreproxy}
Theorem~\ref{thm:optimality}'s optimality condition, Eq.~\eqref{eq:optcond}, is a \emph{pointwise}
statement about where the true density equals a constant at the region boundary; it is not, in
general, equivalent to minimizing the \emph{global} discrepancy
$\mathbb E_X[\mathrm{KL}(p(\cdot\mid X)\|f_{c,\rho}(\cdot\mid X))]$ that
Proposition~\ref{prop:logscore-kl} shows the log score targets, since a density can be globally close
in KL divergence while still being systematically wrong exactly at the region's boundary, or vice
versa. Selecting $(c,\rho)$ by the log score is therefore a well-motivated, strictly-proper-scoring-rule
proxy for solving Eq.~\eqref{eq:optcond} -- not a construction that provably achieves it. This is the
precise sense in which MELD-HPD's shape selection is theoretically well grounded but, unlike
Proposition~\ref{prop:validity}'s coverage guarantee, not an exact guarantee of narrower width.
\end{remark}

\subsection{Selecting $\theta$, and a closed-form diagnostic for $\rho$}
$\theta$ does not enter $f_{c,\rho}$ (Eq.~\eqref{eq:fcrho}); it only tilts the threshold, so
$L(c,\rho)$ of Eq.~\eqref{eq:logscore} is constant in $\theta$ and cannot rank it. $\theta^\star$ is
therefore still selected by $\arg\min$ tune-fold width $\bar W(c^\star,\rho^\star,\theta)$ at the
now-fixed $(c^\star,\rho^\star)$ -- its own proper objective, and a single well-behaved dimension far
less prone to the overfitting risk of a joint $(c,\rho,\theta)$ width $\arg\min$.

As an interpretable anchor -- a diagnostic, not part of the selection rule -- MELD-HPD also reports a
closed-form least-squares estimate of how much epistemic spread the residuals support.

\begin{lemma}[Least-squares epistemic-reliability estimate]
\label{lem:closedrho}
Fix $c$, write $A_i=A(\bm x_i;c)$, $E_i=\sum_k w_k(\bm x_i)(\mu_k(\bm x_i)-\hat y(\bm x_i))^2$ (the
\emph{raw}, $\rho=1$, epistemic term), and $R_i^2=(y_i-\hat y(\bm x_i))^2$. The minimizer over
$\rho\in\mathbb R$ of $J(\rho)=\sum_{i=1}^n (A_i+\rho E_i - R_i^2)^2$ is
\begin{equation}
\hat\rho = \frac{\sum_i E_i(R_i^2-A_i)}{\sum_i E_i^2}, \qquad \sum_i E_i^2>0,
\label{eq:closedrho}
\end{equation}
and the reported diagnostic is $\max(\hat\rho,0)$, the projection of Eq.~\eqref{eq:closedrho} onto
the domain $\rho\ge0$ that Eq.~\eqref{eq:etamu}'s $\sqrt\rho$ requires. $\hat\rho$ is undefined
(reported as \texttt{NaN}) when $\sum_i E_i^2=0$, i.e.\ whenever $K=1$, where the epistemic term
carries no information for any $\rho$ to be fit against.
\end{lemma}
\begin{proof}
$J$ is quadratic and convex in $\rho$ (since $J''(\rho)=2\sum_i E_i^2\ge0$). Setting
$J'(\rho)=2\sum_i E_i(A_i+\rho E_i - R_i^2)=0$ and solving for $\rho$ gives
$\rho\sum_i E_i^2=\sum_i E_i(R_i^2-A_i)$, i.e.\ Eq.~\eqref{eq:closedrho}; convexity makes this
stationary point the global minimizer whenever $\sum_i E_i^2>0$.
\end{proof}

Equation~\eqref{eq:closedrho} is exactly the ordinary-least-squares fit, in the single covariate
$E_i$ with fixed offset $A_i$, of the reshaped law-of-total-variance identity
$\sigma_{\mathrm{tot}}^2(\bm x;c,\rho)=A(\bm x;c)+\rho E(\bm x)$ to the observed squared residuals --
``how much of the model's reported epistemic spread the residuals actually support'' at the chosen
$c^\star$.

\subsection{Boundary recovery and validity}
The set of reachable $(c,\rho,\theta)$ configurations is identical to SETA-HPD's,
$\mathcal C\times\mathcal R\times\Theta$, and still contains $(1,1,0)$ (plain HPD calibration) and
$(c,1,\theta)$ for every $(c,\theta)$ (STA-HPD), by the same substitutions as
Proposition~\ref{prop:setarecovery}.

\begin{remark}[Same grid, weaker floor]
\label{rem:meldweakfloor}
Because $(c^\star,\rho^\star)$ is chosen to maximize $L(c,\rho)$ rather than minimize
$\bar W(c,\rho,\theta)$, and $(1,1)\in\mathcal C\times\mathcal R$, MELD-HPD's tuning objective
satisfies $L(c^\star,\rho^\star)\ge L(1,1)$ -- a floor on \emph{log density}, not on width. This does
\emph{not} imply $\bar W(c^\star,\rho^\star,\theta^\star)\le\bar W(1,1,0)$: unlike
Propositions~\ref{prop:setarecovery} and \ref{prop:easerecovery}'s width-argmin chains, MELD-HPD
carries no guarantee that its tune-fold width is at most plain HPD's or STA-HPD's, only that
$\theta^\star$ is itself width-optimal \emph{given} the log-score-selected $(c^\star,\rho^\star)$,
since $\theta=0\in\Theta$ is still on the one grid that $\theta^\star$ is chosen from. The worked
comparison below shows this gap can be realized in practice, not merely in principle.
\end{remark}

Validity is unaffected by any of this: the coverage guarantee of Proposition~\ref{prop:validity}
depends only on $(c,\rho,\theta)$ being fixed by a rule that is a function of the tune fold alone,
independent of the calib-fold labels supplying $\hat t$ -- it does not depend on \emph{which}
objective that rule optimizes. Since $(c^\star,\rho^\star)$ is selected from $\mathcal D_{\mathrm{tune}}$
by Eq.~\eqref{eq:logscore} and $\theta^\star$ from the same fold by $\bar W$, both disjoint from
$\mathcal D_{\mathrm{calib}}$, the resulting score is a frozen, label-independent function of
$(\bm x_i,y_i)$ by the time it is conformalized, exactly as in Proposition~\ref{prop:validity}.

\subsection{Fitting}
Fitting replaces Algorithm~\ref{alg:stahpd}'s single width-$\arg\min$ pass over
$\mathcal C\times\mathcal R\times\Theta$ with two stages per split: (1) $(c^\star,\rho^\star)$ is the
majority vote, over the $n_{\mathrm{repeats}}$ splits, of the per-split $\arg\max_{c,\rho} L(c,\rho)$
on that split's tune fold (each per-point log-density floored at $-10^3$ nats before averaging, so a
single underflow outlier cannot send a cell to $-\infty$); (2) at the now-fixed $(c^\star,\rho^\star)$,
$\theta^\star$ is the majority vote of the per-split $\arg\min_\theta \bar W(c^\star,\rho^\star,\theta)$.
The final threshold is the median-of-calib-fold-quantiles construction of Eq.~\eqref{eq:medianquantile},
unchanged.

\subsection{A worked comparison across the whole family}
% Why this table is placed here, after both new sections, rather than split across them: it needs
% both EASE-HPD's and MELD-HPD's numbers to make its point (that the width floor holds for one
% extension and visibly breaks for the other, on the identical config and split construction), so it
% only says something once both are on the page.
On the identical checkpoint, split construction, and nominal coverage as the EASE-HPD worked example,
Table~\ref{tab:allfive} reports all five members of the family discussed in this paper.

\begin{table}[H]
\centering
\caption{All five HPD-family calibrators on one held-out MOG checkpoint (Synthetic, single-expert
probabilistic head, seed $4022$, $1-\alpha=90\%$), identical tune/calib construction throughout.}
\label{tab:allfive}
\begin{tabular}{@{}lrrl@{}}
\toprule
\textbf{Method} & \textbf{Coverage} & \textbf{Width} & \textbf{Selected parameters} \\
\midrule
MoG-HPD & $0.9087$ & $1.5757$ & --- \\
STA-HPD & $0.8933$ & $1.3473$ & $(c,\theta){=}(0.8,1.5)$ \\
SETA-HPD & $0.8933$ & $1.3473$ & $(c,\rho,\theta){=}(0.8,0,1.5)$ \\
EASE-HPD & $0.8933$ & $1.3473$ & $(c,\rho,\theta,\phi){=}(0.8,0,1.5,{-1.0})$ \\
MELD-HPD & $0.8940$ & $1.3708$ & $(c,\rho,\theta){=}(1.3,0,0.75)$ \\
\bottomrule
\end{tabular}
\end{table}

STA-HPD, SETA-HPD, and EASE-HPD agree to four decimals, as their respective floor propositions
guarantee on this config once $\rho^\star=0$ and $\phi$'s inertness are accounted for
(Remark~\ref{rem:phiinert}). MELD-HPD is the one visible exception, at $1.3708$ -- $1.7\%$
\emph{wider} than the other four -- and the logged tuning surfaces show exactly why. The log-score
surface is monotone increasing in $c$ on this config's grid ($L(0.8,0)=-0.4528$, $L(1,0)=-0.3885$,
$L(1.3,0)=-0.3634$, all at $\rho=0$), so MELD-HPD's proper-scoring-rule selector correctly picks the
largest available $c$, $c^\star=1.3$. But the tune-fold \emph{width} surface at that same cell,
$\bar W(1.3,0,0.75)=1.4705$, is not the grid's minimum; that minimum, $\bar W(0.8,0,1.5)=1.4402$, sits
at the smaller $c=0.8$ that STA-HPD, SETA-HPD, and EASE-HPD all select instead. This is
Remark~\ref{rem:logscoreproxy} realized on real data: the density-closest member of the family, by
the strictly proper log score, is demonstrably not the width-narrowest one on this particular config,
exactly the gap a KL-motivated selector does not, and by Remark~\ref{rem:meldweakfloor}'s argument
structurally cannot, close.

\section{Conclusion}
\label{sec:conclusion}
% Why the conclusion restates the two knobs and their theoretical grounding
% before the empirical summary, rather than leading with the numbers: a
% conference-paper conclusion should let a reader who skipped straight from
% the abstract to the conclusion still recover the paper's core argument
% (what's broken, what fixes it, why the fix is guaranteed not to hurt) --
% the empirical numbers are supporting evidence for that argument, not the
% argument itself.
STA-HPD generalizes highest-density-region conformal calibration for mixture-of-experts regressors
along exactly the two axes that a fixed-density level set cannot address on its own: the shape of
the model's own heteroscedasticity and the level at which that heteroscedasticity is thresholded.
Both knobs are selected by a conformal-risk-minimizing grid search that provably contains Standard
CP, CP-VS (in the single-component case), and plain HPD calibration as boundary cases, and the
entire construction retains exact finite-sample marginal coverage for any fixed choice of
$(c,\theta)$ regardless of model correctness, since validity and efficiency are cleanly separated
throughout. Empirically the method yields consistent, never-negative width reductions relative to
plain HPD calibration, with the largest gains concentrated where the fitted mixture is most
genuinely multimodal. Sections~\ref{sec:seta}--\ref{sec:meld} extend the same construction along
three further axes, each isolating a different piece of the MoG variance decomposition or a
different way of selecting it: SETA-HPD reshapes the between-component (epistemic) variance that
STA-HPD leaves untouched; EASE-HPD tilts the threshold by the point's epistemic \emph{share} rather
than only its total scale; and MELD-HPD replaces tune-fold width $\arg\min$ with a strictly proper
scoring rule for the two density-shaping exponents, trading the width-argmin family's tuning-fold
floor guarantee for a population-level KL-optimality argument instead
(Remarks~\ref{rem:logscoreproxy}--\ref{rem:meldweakfloor}). Every extension is validated by the same
exchangeability argument and, where a width-argmin selector is used, provably contains its ancestors
as reachable grid points; MELD-HPD's worked comparison (Table~\ref{tab:allfive}) shows this design
space is genuine, not merely nominal, since the two selection principles can and do disagree on real
data. Natural extensions include a per-expert (rather than scalar) shape exponent,
adaptive grid refinement in place of a fixed $(c,\theta)$ lattice, and extending the threshold-field
tilt to depend on additional covariates of the routing distribution beyond the total scale and
epistemic share.

% Why these ten references specifically: each was chosen to back a
% distinct claim made in the text, not as a generic literature survey --
% vovk2005algorithmic/lei2018distribution for the split-conformal machinery
% eq:qlevel relies on; romano2019conformalized for CQR (the one baseline
% whose region shape differs qualitatively, discussed in Sec. III);
% boxtiao1973bayesian/hyndman1996computing for the classical, pre-conformal
% highest-density-region background in the new Sec. II-B;
% izbicki2020cd/sadinle2019least for the HPD/level-set construction
% Theorem 1 formalizes in the conformal setting; gibbs2021adaptive/
% tibshirani2019conformal for the adaptive/weighted-conformal stabilization
% ideas the repeated-splitting procedure in Sec. II-E is modeled after;
% jordan1994hierarchical for the MoE architecture itself, since the paper
% assumes the reader knows what a gated mixture of experts is.
\begin{thebibliography}{10}

\bibitem{vovk2005algorithmic}
V.~Vovk, A.~Gammerman, and G.~Shafer, \emph{Algorithmic Learning in a Random World}. Springer, 2005.

\bibitem{lei2018distribution}
J.~Lei, M.~G'Sell, A.~Rinaldo, R.~J.~Tibshirani, and L.~Wasserman, ``Distribution-free predictive
inference for regression,'' \emph{J. Amer. Statist. Assoc.}, vol.~113, no.~523, pp. 1094--1111,
2018.

\bibitem{romano2019conformalized}
Y.~Romano, E.~Patterson, and E.~Cand\`es, ``Conformalized quantile regression,'' in \emph{Advances
in Neural Information Processing Systems (NeurIPS)}, 2019.

\bibitem{izbicki2020cd}
R.~Izbicki, G.~Shimizu, and R.~B.~Stern, ``Flexible distribution-free conditional predictive bands
using density estimators,'' in \emph{Proc. International Conference on Artificial Intelligence and
Statistics (AISTATS)}, 2020.

\bibitem{sadinle2019least}
M.~Sadinle, J.~Lei, and L.~Wasserman, ``Least ambiguous set-valued classifiers with bounded error
levels,'' \emph{J. Amer. Statist. Assoc.}, vol.~114, no.~525, pp. 223--234, 2019.

\bibitem{gibbs2021adaptive}
I.~Gibbs and E.~Cand\`es, ``Adaptive conformal inference under distribution shift,'' in
\emph{Advances in Neural Information Processing Systems (NeurIPS)}, 2021.

\bibitem{tibshirani2019conformal}
R.~J.~Tibshirani, R.~Foygel~Barber, E.~Cand\`es, and A.~Ramdas, ``Conformal prediction under
covariate shift,'' in \emph{Advances in Neural Information Processing Systems (NeurIPS)}, 2019.

\bibitem{jordan1994hierarchical}
M.~I.~Jordan and R.~A.~Jacobs, ``Hierarchical mixtures of experts and the EM algorithm,''
\emph{Neural Computation}, vol.~6, no.~2, pp. 181--214, 1994.

\bibitem{boxtiao1973bayesian}
G.~E.~P.~Box and G.~C.~Tiao, \emph{Bayesian Inference in Statistical Analysis}. Addison-Wesley,
1973.

\bibitem{hyndman1996computing}
R.~J.~Hyndman, ``Computing and graphing highest density regions,'' \emph{The American
Statistician}, vol.~50, no.~2, pp. 120--126, 1996.

\bibitem{gneiting2007strictly}
T.~Gneiting and A.~E.~Raftery, ``Strictly proper scoring rules, prediction, and estimation,''
\emph{J. Amer. Statist. Assoc.}, vol.~102, no.~477, pp. 359--378, 2007.

% Added for the Related Work section: the three lines of prior work whose
% *objective* (post-hoc width reduction) or *mechanism* (normalized scores)
% overlaps STA-HPD closely enough to require explicit positioning.
% papadopoulos2011regression: the normalized/locally-adaptive nonconformity
% score STA-HPD's theta=1 axis reduces to. xie2024boosted: the closest prior
% work (post-hoc score reshaping for width). kiyani2024length: the
% length-optimization-under-validity framework STA-HPD is a specialized point
% within.
\bibitem{papadopoulos2011regression}
H.~Papadopoulos, V.~Vovk, and A.~Gammerman, ``Regression conformal prediction with nearest
neighbours,'' \emph{J. Artificial Intelligence Research}, vol.~40, pp. 815--840, 2011.

\bibitem{xie2024boosted}
R.~Xie, R.~Foygel~Barber, and E.~J.~Cand\`es, ``Boosted conformal prediction intervals,'' in
\emph{Advances in Neural Information Processing Systems (NeurIPS)}, 2024.

\bibitem{kiyani2024length}
S.~Kiyani, G.~Pappas, and H.~Hassani, ``Length optimization in conformal prediction,'' in
\emph{Advances in Neural Information Processing Systems (NeurIPS)}, 2024.

\end{thebibliography}

\end{document}
